\documentclass[a4paper,11pt]{article}
\usepackage[utf8]{inputenc}
\usepackage[T2A]{fontenc}
\usepackage[english,russian]{babel}
\usepackage[margin=2.2cm]{geometry}
\usepackage{amsmath,amssymb,amsthm}
\usepackage{stmaryrd}
\usepackage{proof}
\usepackage{tikz}
\usepackage{enumitem}

\theoremstyle{definition}
\newtheorem*{thm}{Теорема}
\newtheorem*{lmm}{Лемма}
\newtheorem*{axm}{Аксиома}
\newtheorem*{dfn}{Определение}

\newcommand{\sem}[2]{\llbracket #1 \rrbracket_{#2}}
\newcommand{\proves}{\vdash}
\newcommand{\True}{\text{И}}
\newcommand{\False}{\text{Л}}
\newcommand{\amp}{\mathbin{\&}}
\newcommand{\cn}[1]{\langle\!\!\langle #1 \rangle\!\!\rangle}
\newcommand{\gn}[1]{\ulcorner #1 \urcorner}
\newcommand{\Sh}[1]{S\langle #1\rangle}
\newcommand{\Rr}[1]{R\langle #1\rangle}
\newcommand{\Mn}[1]{M\langle #1\rangle}
\newcommand{\todo}[1]{\textbf{TODO:} #1}
\newcommand{\upb}{\operatorname{upb}}
\newcommand{\rk}{\operatorname{rk}}

\title{Математическая логика}
\author{Автор: Вячеслав Чепелин\\[2pt] Лектор: Штукенберг Д.Г.}
\date{}

\begin{document}
\maketitle
\tableofcontents

\section{Исчисление высказываний}
\begin{dfn}
Предметный язык~--- формализация мат.\ утверждений.
\end{dfn}
\begin{dfn}
Метаязык (или язык исследователя)~--- язык, которым пользуются для формирования
математических утверждений.
\end{dfn}
\begin{dfn}
Высказывание~--- строка, которая либо атомарное высказывание (то есть какая-то
переменная), либо составное (составлено с помощью отрицания, дизъюнкции,
конъюнкции или импликации из двух высказываний).
\end{dfn}
\begin{dfn}
Схема высказываний~--- высказывание, где мы можем использовать метапеременные
(вместо которых можно подставлять что угодно).
\end{dfn}
\begin{dfn}
Язык исчисления высказываний~--- язык, порождаемый грамматикой высказываний.
\end{dfn}
\begin{dfn}
Оценка высказываний~--- множество истинностных значений (множество $V =
\{\True, \False\}$) и функция оценки переменных. Все оценки составных
высказываний рекурсивно задаются.
\end{dfn}
\begin{dfn}
$\alpha$~--- тавтология (общезначима), если $\models \alpha$ или по-другому, что
$\alpha$ истинна при всех оценках.
\end{dfn}
\begin{dfn}
$\gamma_1, \dots, \gamma_n \models \alpha$~--- следствие, то есть если данные
$\gamma$ выполнены, то $\alpha$ выполнено.
\end{dfn}
\begin{dfn}
Доказательством назовем конечную последовательность $\delta_1, \dots, \delta_n$,
где каждое $\delta_i$ либо аксиома, либо получается по правилу Modus Ponens
($\alpha, \alpha \to \beta \Rightarrow \beta$).
\end{dfn}
В классической логике есть 10 аксиом, которые задаются схемами. Назовем эти
аксиомы аксиомами исчисления высказываний.
\begin{dfn}
Доказательство формулы $\alpha$~--- такое доказательство, что $\delta_n \equiv
\alpha$. Тогда $\alpha$ называется доказуемой.
\end{dfn}
\begin{dfn}
Вывод формулы из гипотез $\gamma_1, \dots, \gamma_n$~--- в доказательстве можем
использовать еще и гипотезы. Обозначается $\gamma_1, \dots, \gamma_n \proves
\alpha$.
\end{dfn}
\begin{dfn}
Теория корректна, если любое доказуемое утверждение общезначимо.
\end{dfn}
\begin{dfn}
Теория полна, если любое общезначимое утверждение доказуемо.
\end{dfn}
\begin{dfn}
Теория противоречива, если существует $\alpha$, такое что $\proves \alpha \amp
\neg\alpha$. Эквивалентно, что любая формула доказуема, или для некоторой
$\alpha$ имеет место $\proves \alpha$ и $\proves \neg\alpha$.
\end{dfn}

\begin{thm}[о дедукции]
$\Gamma, \alpha \proves \beta \Leftrightarrow \Gamma \proves \alpha \to \beta$.
\end{thm}
\begin{proof}
Доказательство в левую сторону очевидно. Докажем, что из $\Gamma, \alpha \proves
\beta$ следует $\Gamma \proves \alpha \to \beta$.

Пусть $\Gamma, \alpha \proves \beta$, т.е.\ существует вывод $D = (\delta_1,
\dots, \delta_m)$, где $\delta_m = \beta$. Построим вывод $\alpha \to \beta$ из
$\Gamma$ индукцией по длине вывода $D$.

\emph{База:} очевидно.

\emph{Шаг индукции:} Предположим, для всех выводов длины $\le m$ утверждение
верно. Рассмотрим вывод $D$ длины $m+1$:
\[ \delta_1, \dots, \delta_m, \delta_{m+1} (= \beta) \]
По индукционному предположению, уже построены выводы $\zeta_i$ из $\Gamma$, такие
что $\zeta_i = \alpha \to \delta_i$ для $i = 1, \dots, m$.

Разберём возможные случаи для $\delta_{m+1}$:
\begin{enumerate}
\item $\delta_{m+1}$~--- аксиома или $\delta_{m+1} \in \Gamma$: Тогда $\Gamma
  \proves \delta_{m+1}$. Добавим к построенным $\zeta_1, \dots, \zeta_m$ шаги:
  \begin{align*}
  &\delta_{m+1} \to (\alpha \to \delta_{m+1}) &&\text{(акс.\ 1)}\\
  &\delta_{m+1}\\
  &\alpha \to \delta_{m+1} &&\text{(MP)}
  \end{align*}
\item $\delta_{m+1} \equiv \alpha$, $\alpha \to \alpha$ выводимо всегда.
\item $\delta_{m+1}$ получен по MP из $\delta_j$ и $\delta_k$ ($\delta_k \equiv
  \delta_j \to \delta_{m+1}$), где $j, k < m+1$: По индукционному предположению
  уже имеем $\Gamma \proves \alpha \to \delta_j$ и $\Gamma \proves \alpha \to
  (\delta_j \to \delta_{m+1})$. По акс.\ 2:
  \[ (\alpha \to \delta_j) \to (\alpha \to (\delta_j \to \delta_{m+1})) \to (\alpha \to \delta_{m+1}) \]
  Применив MP дважды, получаем $\Gamma \proves \alpha \to \delta_{m+1}$.
\end{enumerate}
Таким образом, во всех случаях $\Gamma \proves \alpha \to \delta_{m+1}$, т.е.
$\Gamma \proves \alpha \to \beta$.
\end{proof}

\section{Теорема о полноте исчисления высказываний}
\begin{dfn}[Условное отрицание]
Зададим некоторую оценку переменных, что $\llbracket a\rrbracket = x$. Тогда
условным отрицанием формулы $\alpha$ назовем следующую формулу $\cn{\alpha}$:
\[
\cn{\alpha} =
\begin{cases}
\alpha, & x = \True\\
\neg\alpha, & x = \False
\end{cases}
\]
Также, если $\Gamma := \gamma_1, \gamma_2, \dots, \gamma_n$, то за $\cn{\Gamma}$
обозначим $\cn{\gamma_1}, \cn{\gamma_2}, \dots \cn{\gamma_n}$.
\end{dfn}

\begin{thm}[о полноте]
Если $\models \alpha$, то $\proves \alpha$.
\end{thm}
\begin{proof}
\begin{enumerate}
\item Построим таблицы истинности для каждой связки ($\star$) и докажем в них
  каждую строку: $\cn{\varphi}, \cn{\psi} \proves \cn{\varphi \star \psi}$.
\item Построим таблицу истинности для $\alpha$ и докажем в ней каждую строку:
  $\cn{\Xi} \proves \cn{\alpha}$.
\item Если формула общезначима, то в ней все строки будут иметь вид $\cn{\Xi}
  \proves \alpha$, потому от гипотез мы сможем избавиться и получить требуемое
  $\proves \alpha$.
\end{enumerate}
Приступим:

\emph{Шаг 1. Лемма о связках.} Запись $\cn{\varphi}, \cn{\psi} \proves
\cn{\varphi \star \psi}$ сводится к 14 утверждениям:
\begin{itemize}
\item $\neg\varphi, \neg\psi \proves \neg(\varphi \amp \psi)$ и $\neg\varphi,
  \neg\psi \proves (\varphi \to \psi)$
\item $\neg\varphi, \psi \proves \neg(\varphi \amp \psi)$ и $\neg\varphi, \psi
  \proves (\varphi \to \psi)$
\item $\varphi, \neg\psi \proves \neg(\varphi \amp \psi)$ и $\varphi, \neg\psi
  \proves \neg(\varphi \to \psi)$
\item $\varphi, \psi \proves (\varphi \amp \psi)$ и $\varphi, \psi \proves
  (\varphi \to \psi)$
\item $\neg\varphi, \neg\psi \proves \neg(\varphi \vee \psi)$ и $\varphi \proves
  \neg\neg\varphi$
\item $\neg\varphi, \psi \proves (\varphi \vee \psi)$ и $\neg\varphi \proves
  \neg\varphi$
\item $\varphi, \neg\psi \proves (\varphi \vee \psi)$
\item $\varphi, \psi \proves (\varphi \vee \psi)$
\end{itemize}

\emph{Шаг 2. Обобщение на любую формулу.}

\begin{lmm}[Условное отрицание формул]
Пусть пропозициональные переменные $\Xi := \{X_1, \dots, X_n\}$~--- все
переменные, которые используются в формуле $\alpha$. И пусть задана некоторая
оценка переменных. Тогда $\cn{\Xi} \proves \cn{\alpha}$.
\end{lmm}
\begin{proof}
Индукция по длине формулы $\alpha$.

\emph{База:} Формула $\alpha$~--- атомарная, т.е.\ $\alpha \equiv X_i$. Тогда
при любом $\Xi$ выполнено $\cn{\Xi}_{X_i := \True} \proves X_i$ и
$\cn{\Xi}_{X_i := \False} \proves \neg X_i$.

\emph{Индукционный переход:} $\alpha \equiv \varphi \star \psi$, причём
$\cn{\Xi} \proves \cn{\varphi}$ и $\cn{\Xi} \proves \cn{\psi}$. Тогда построим
вывод:
\begin{align*}
&(1)\dots(n)\ \cn{\varphi} &&\text{(индукционное предположение)}\\
&(n+1)\dots(k)\ \cn{\psi} &&\text{(индукционное предположение)}\\
&(k+1)\dots(l)\ \cn{\varphi \star \psi} &&\text{(лемма о связках: $\cn{\varphi}$ и $\cn{\psi}$ доказаны выше,}\\
&&&\text{значит, их можно использовать как гипотезы)}
\end{align*}
Откуда лемма доказана.
\end{proof}

\emph{Шаг 3. Избавляемся от гипотез.}

\begin{lmm}[об устранении посылок]
Пусть при всех оценках переменных $\cn{\Xi} \proves \alpha$, тогда $\proves
\alpha$.
\end{lmm}
\begin{proof}
Индукция по количеству переменных $n$.

\emph{База:} $n = 0$. Тогда $\proves \alpha$ есть из условия.

\emph{Индукционный переход:} Пусть $\cn{X_1, X_2, \dots X_{n+1}} \proves
\alpha$. Рассмотрим выводы:
\[
\cn{X_1, X_2, \dots X_n}, X_{n+1} \proves \alpha
\qquad
\cn{X_1, X_2, \dots X_n}, \neg X_{n+1} \proves \alpha
\]
По лемме об исключении допущения тогда $\cn{X_1, X_2, \dots X_n} \proves \alpha$.
(Лемма об исключении допущения: $\alpha \proves \beta$ и $\neg\alpha \proves
\beta$, то $\proves \beta$. Доказывается очень просто с помощью теоремы о
дедукции.)

При этом, $\cn{X_1, X_2, \dots X_n} \proves \alpha$ при всех оценках переменных
$X_1, \dots X_n$. Значит, $\proves \alpha$ по индукционному предположению.
\end{proof}
\end{proof}

\section{Топологические пространства}
\begin{dfn}
Упорядоченная пара $\langle X, \Omega\rangle$ где $X$~--- множество, $\Omega
\subseteq \mathcal{P}(X)$ (называется топологией), удовлетворяющая:
\begin{enumerate}
\item $\varnothing \in \Omega$, $X \in \Omega$.
\item Конечное пересечение элементов $\Omega$ принадлежит $\Omega$.
\item Объединение любого семейства элементов $\Omega$ принадлежит $\Omega$.
\end{enumerate}
\end{dfn}
\begin{dfn}
Открытые множества~--- элементы топологии $\Omega$.
\end{dfn}
\begin{dfn}
Замкнутое множество~--- множество, дополнение до $X$ которого является открытым.
\end{dfn}
\begin{dfn}
Внутренность множества ($A^\circ$)~--- наибольшее открытое множество,
содержащееся в $A$.
\end{dfn}

\subsection*{Примеры топологий}
\begin{itemize}
\item Топология стрелки: на $\mathbb{R}$, открытыми являются множества вида
  $[a, +\infty)$ и пустое множество.
\item Топология Зарисского: на $\mathbb{R}$, открытыми являются $\mathbb{R}$,
  $\varnothing$ и все множества, дополнения которых конечны.
\item Топология на дереве (лесе): Пусть задан лес с множеством вершин $V$ и
  отношением порядка $\preceq$ (предок-потомок: $a \preceq b$, если $b$~---
  потомок $a$). Открытое множество вместе с каждой вершиной содержит всех её
  потомков.
\end{itemize}

\begin{dfn}
Непрерывная функция (между топологическими пространствами)~--- прообраз любого
открытого множества в $Y$ является открытым множеством в $X$.
\end{dfn}
\begin{dfn}
Компактное множество~--- множество, из любого его открытого покрытия можно
выбрать конечное подпокрытие.
\end{dfn}
\begin{dfn}
Подпространство топологического пространства~--- пространство $\langle X_1,
\Omega_1\rangle$ называется подпространством $\langle X, \Omega\rangle$, если
$X_1 \subseteq X$ и $\Omega_1 = \{A \cap X_1 \mid A \in \Omega\}$.
\end{dfn}
\begin{dfn}
Связное пространство~--- не существует таких двух непустых открытых множеств
$A, B \in \Omega$, что $A \cup B = X$ и $A \cap B = \varnothing$.
\end{dfn}
\begin{dfn}
Связное множество~--- если соответствующее ему подпространство связно.
\end{dfn}

\section{Интуиционистская логика, вывод, интерпретация, решетки}
Основные положения интуиционизма:
\begin{itemize}
\item Математика не формальна.
\item Математика независима от окружающего мира.
\item Логика зависит от математики, а не наоборот.
\end{itemize}

\begin{dfn}
Доказательства чистого существования (конструктивный подход): Объект считается
существующим, если для него предоставлен эффективный метод построения, а не
только доказательство невозможности его отсутствия.
\end{dfn}
\begin{dfn}[BHK-интерпретация (Брауэр, Гейтинг, Колмогоров)]
Конструктивная интерпретация логических связок, где истинность высказывания
означает наличие его построения (доказательства).
\begin{itemize}
\item $\alpha \wedge \beta$ построено, если построены $\alpha$ и $\beta$.
\item $\alpha \vee \beta$ построено, если построено $\alpha$ или $\beta$, и
  известно, какое именно.
\item $\alpha \to \beta$ построено, если имеется метод, преобразующий любое
  построение $\alpha$ в построение $\beta$.
\item $\bot$~--- конструкция, не имеющая построения.
\item $\neg\alpha$ построено, если построено $\alpha \to \bot$.
\end{itemize}
\end{dfn}
\begin{dfn}
Интуиционистское исчисление высказываний (ИИВ) в гильбертовском стиле:
Получается из аксиоматики классического исчисления высказываний (КИВ) заменой
аксиомы снятия двойного отрицания ($\neg\neg\alpha \to \alpha$) на аксиому
($\alpha \to (\neg\alpha \to \beta)$) (принцип взрыва).
\end{dfn}
\begin{dfn}
Интуиционистское исчисление высказываний (натуральный вывод):
\begin{itemize}
\item Секвенция: Выражение вида $\Gamma \proves \alpha$, где $\Gamma$~---
  множество гипотез, $\alpha$~--- заключение.
\item Аксиома: $\Gamma, \alpha \proves \alpha$.
\item Правила введения связок:
  \begin{itemize}
  \item Если $\Gamma, \alpha \proves \beta$, то $\Gamma \proves \alpha \to \beta$.
  \item Если $\Gamma \proves \alpha$, то $\Gamma \proves \alpha \vee \beta$. Если
    $\Gamma \proves \beta$, то $\Gamma \proves \alpha \vee \beta$.
  \item Если $\Gamma \proves \alpha$ и $\Gamma \proves \beta$, то $\Gamma
    \proves \alpha \wedge \beta$.
  \end{itemize}
\item Правила удаления связок:
  \begin{itemize}
  \item Если $\Gamma \proves \alpha$ и $\Gamma \proves \alpha \to \beta$, то
    $\Gamma \proves \beta$.
  \item Если $\Gamma \proves \alpha \vee \beta$, $\Gamma \proves \alpha \to
    \gamma$ и $\Gamma \proves \beta \to \gamma$, то $\Gamma \proves \gamma$.
  \item Если $\Gamma \proves \alpha \wedge \beta$, то $\Gamma \proves \alpha$.
    Если $\Gamma \proves \alpha \wedge \beta$, то $\Gamma \proves \beta$.
  \item Если $\Gamma \proves \bot$, то $\Gamma \proves \alpha$ (из лжи следует
    всё).
  \end{itemize}
\end{itemize}
Его мы записываем в виде дерева.
\end{dfn}
\begin{dfn}[Закон исключённого третьего ($\alpha \vee \neg\alpha$)]
В BHK-интерпретации не является общезначимым, так как для нерешённой проблемы у
нас нет построения ни для $\alpha$, ни для $\neg\alpha$.
\end{dfn}
\begin{dfn}[Изоморфизм Карри--Ховарда]
Прямая аналогия между логикой и теорией типов в программировании.
\begin{itemize}
\item Программа ($\lambda$-выражение) $\leftrightarrow$ Доказательство.
\item Тип выражения $\leftrightarrow$ Высказывание.
\item Тип функции ($A \to B$) $\leftrightarrow$ Импликация ($A \to B$).
\item Упорядоченная пара ($A \times B$) $\leftrightarrow$ Конъюнкция ($A \wedge B$).
\item Алгебраический тип ($A + B$) $\leftrightarrow$ Дизъюнкция ($A \vee B$).
\item Необитаемый тип (\texttt{void}, $\bot$) $\leftrightarrow$ Ложь ($\bot$).
\end{itemize}
\end{dfn}
\begin{dfn}[Решётка]
Упорядоченная пара $\langle X, \preceq\rangle$, где $X$~--- множество, $\preceq$~---
частичный порядок на $X$, такой что для любых $a, b \in X$ определены:
\begin{itemize}
\item Сумма (супремум): $a + b = \sup\{a, b\}$~--- наименьший элемент $c$ такой,
  что $a \preceq c$ и $b \preceq c$.
\item Произведение (инфимум): $a \cdot b = \inf\{a, b\}$~--- наибольший элемент
  $c$ такой, что $c \preceq a$ и $c \preceq b$.
\end{itemize}
\end{dfn}
\begin{dfn}[Дистрибутивная решётка]
Решётка, в которой для любых $a, b, c$ выполняется: $a \cdot (b + c) = (a \cdot
b) + (a \cdot c)$.
\end{dfn}
\begin{dfn}[Псевдодополнение]
В решётке псевдодополнением $a \to b$ называется наибольший элемент из множества
$\{x \mid a \cdot x \preceq b\}$.
\end{dfn}
\begin{dfn}[Импликативная решётка]
Решётка, в которой для любых двух элементов существует псевдодополнение.
\end{dfn}
$0$ и $1$ в решётке: $0$~--- наименьший элемент решётки, $1$~--- наибольший
элемент решётки.
\begin{dfn}[Псевдобулева алгебра (алгебра Гейтинга)]
Импликативная решётка с наименьшим элементом $0$. В ней определяется операция
отрицания: $\neg a = a \to 0$.
\end{dfn}
\begin{dfn}[Булева алгебра]
Псевдобулева алгебра, в которой для любого элемента $a$ выполняется закон
исключённого третьего в алгебраической форме: $a + \neg a = 1$.
\end{dfn}

\section{Продолжение решеток, алгебра Линденбаума, модели Крипке}
\begin{dfn}
Пусть $L$~--- множество всех высказываний. Тогда алгебра Линденбаума $\mathcal{L}
= L/{\simeq}$, где отношение эквивалентности построено, если $a \preceq b$, $b
\preceq a$, где $\preceq$~--- предпорядок: $\alpha \preceq \beta := \alpha
\proves \beta$ в ИИВ.
\end{dfn}
\begin{thm}
$\mathcal{L}$~--- псевдобулева алгебра. Доказательство~--- надо показать, что
$\preceq$ отношение порядка на $\mathcal{L}$.
\end{thm}
\begin{dfn}
Пусть некоторое исчисление оценивается значениями из решетки. Оценка согласована
с исчислением, если «и» это умножение, «или» это сложение и так далее.
\end{dfn}
\begin{thm}
$\llbracket\alpha\rrbracket = [\alpha]_{\mathcal{L}}$.
\end{thm}

\begin{thm}
Интуиционистское исчисление полно в псевдобулевых алгебрах. То есть если
$\models \alpha$ во всех псевдобулевых, то $\proves \alpha$.
\end{thm}
\begin{proof}
Возьмем в качестве модели исчисления алгебру Линденбаума. Пусть $\models
\alpha$, тогда $\llbracket\alpha\rrbracket = 1$ во всех псевдобулевых, в том
числе и $\llbracket\alpha\rrbracket = 1_{\mathcal{L}}$. То есть
$[\alpha]_{\mathcal{L}} = [A \to A]_{\mathcal{L}}$. То есть $A \to A \proves
\alpha$, откуда уже следует искомое.

Идейно мы просто перекидываем все в алгебру Линденбаума и пользуемся ее
свойствами.
\end{proof}

\begin{dfn}[Модель Крипке]
Упорядоченная тройка $\langle\mathcal{W}, \preceq, \Vdash\rangle$, где:
\begin{itemize}
\item $\mathcal{W}$~--- множество миров.
\item $\preceq$~--- нестрогий частичный порядок на $\mathcal{W}$.
\item $\Vdash\ \subseteq \mathcal{W} \times P$~--- отношение вынужденности между
  мирами и пропозициональными переменными, монотонное относительно порядка:
  если $W_i \preceq W_j$ и $W_i \Vdash X$, то $W_j \Vdash X$.
\end{itemize}
\end{dfn}
\begin{dfn}[Вынужденность (доопределение для формул)]
Для любого мира $W \in \mathcal{W}$:
\begin{itemize}
\item $W \Vdash \alpha \wedge \beta \Leftrightarrow W \Vdash \alpha$ и $W
  \Vdash \beta$
\item $W \Vdash \alpha \vee \beta \Leftrightarrow W \Vdash \alpha$ или $W
  \Vdash \beta$.
\item $W \Vdash \alpha \to \beta \Leftrightarrow$ для любого мира $W_1$ такого,
  что $W \preceq W_1$ и $W_1 \Vdash \alpha$, выполнено $W_1 \Vdash \beta$.
\item $W \Vdash \neg\alpha \Leftrightarrow$ ни для какого мира $W_1$ такого,
  что $W \preceq W_1$, не верно $W_1 \Vdash \alpha$.
\end{itemize}
\end{dfn}
\begin{dfn}
Формула $\alpha$ истинна в модели Крипке ($\models \alpha$), если $W \Vdash
\alpha$ для всех миров $W \in \mathcal{W}$.
\end{dfn}
\begin{dfn}
Общезначимость в моделях Крипке ($\Vdash_{\mathcal{K}} \alpha$), если она
истинна во всех возможных моделях Крипке.
\end{dfn}

\begin{thm}
Пусть $\langle\mathcal{W}, \preceq, \Vdash\rangle$~--- некоторая модель Крипке.
Тогда она есть корректная модель ИИВ.
\end{thm}
\begin{proof}
Доказательство для древовидного отношения. Заметим, что $V(\alpha) := \{w \in
\mathcal{W} \mid w \Vdash \alpha\}$ открыто в топологии для деревьев. Значит
положим $V = \{S \mid S \subseteq \mathcal{W} \amp S \text{ --- открыто}\}$ и
$\llbracket\alpha\rrbracket = V(\alpha)$, получим алгебру Гейтинга.
\end{proof}

\begin{dfn}
Система $\langle V, T, f_\wedge, f_\vee, f_\to, f_\neg\rangle$, где:
\begin{itemize}
\item $V$~--- множество значений.
\item $T \in V$~--- выделенное значение «истина».
\item $f_\wedge, f_\vee, f_\to : V \times V \to V$~--- функции для связок.
\item $f_\neg : V \to V$~--- функция для отрицания.
\end{itemize}
Оценка называется табличной, если значение любой формулы вычисляется рекурсивно
через эти функции и оценку переменных $f_P : P \to V$.
\end{dfn}
\begin{dfn}
Если $\proves \alpha$ влечет $\llbracket\alpha\rrbracket = T$ при всех оценках
пропозициональных переменных, то $\mathcal{M}$~--- табличная модель.
\end{dfn}
\begin{thm}
Не существует полной конечной табличной модели для ИИВ. (Просят только
формулировку.)
\end{thm}

\section{Гёделева алгебра, разрешимость ИИВ}
\begin{dfn}[Гёделевизация алгебры Гейтинга ($\Gamma(\mathcal{A})$)]
Для алгебры Гейтинга $\mathcal{A} = \langle A, \preceq\rangle$ операция
гёделевизации строит новую алгебру $\Gamma(\mathcal{A}) = \langle A \cup
\{\omega\}, \preceq'\rangle$, где:
\begin{itemize}
\item $\omega$~--- новый элемент.
\item Отношение порядка $\preceq'$ минимальное, удовлетворяющее:
  \begin{itemize}
  \item $a \preceq' b$, если $a \preceq b$ в $\mathcal{A}$ и $a, b \notin
    \{\omega, 1\}$.
  \item $a \preceq' \omega$, если $a \ne 1$ в $\mathcal{A}$.
  \item $\omega \preceq' 1$.
  \end{itemize}
\end{itemize}
$\Gamma(\mathcal{A})$ является гёделевой алгеброй.
\end{dfn}
\begin{dfn}
Отображение $g : \mathcal{A} \to \mathcal{B}$ между алгебрами Гейтинга
называется гомоморфизмом, если оно сохраняет операции и константы:
\[
g(a \wedge b) = g(a) \wedge g(b),\quad g(a \vee b) = g(a) \vee g(b),\quad
g(a \to b) = g(a) \to g(b),\quad g(0_\mathcal{A}) = 0_\mathcal{B},\quad
g(1_\mathcal{A}) = 1_\mathcal{B}
\]
\end{dfn}

\begin{thm}[Дизъюнктивность ИИВ]
Если $\proves \alpha \vee \beta$, то либо $\proves \alpha$, либо $\proves
\beta$.
\end{thm}
\begin{proof}
Пусть $\proves \alpha \vee \beta$. Тогда
$\llbracket\alpha \vee \beta\rrbracket_{\Gamma(\mathcal{L})} = 1$, так как данная
оценка согласована с ИИВ. Тогда $\llbracket\alpha\rrbracket_{\Gamma(\mathcal{L})}
= 1$ или $\llbracket\beta\rrbracket_{\Gamma(\mathcal{L})} = 1$, так как
$\Gamma(\mathcal{L})$~--- гёделева. Пусть
$\llbracket\alpha\rrbracket_{\Gamma(\mathcal{L})} = 1$, тогда из гомоморфизма
$\llbracket\alpha\rrbracket_{\mathcal{L}} = 1$, откуда из полноты $\mathcal{L}$
получаем, что $\proves \alpha$.
\end{proof}

\begin{dfn}
Решётка $\mathcal{A}' = \langle A', \preceq'\rangle$ называется подрешёткой
решётки $\mathcal{A} = \langle A, \preceq\rangle$, если:
\begin{itemize}
\item $A' \subseteq A$.
\item $\preceq' = \preceq \cap (A' \times A')$ (ограничение порядка).
\item Для любых $a, b \in A'$ их супремум и инфимум в $\mathcal{A}'$ совпадают с
  супремумом и инфимумом в $\mathcal{A}$.
\end{itemize}
\end{dfn}

\begin{lmm}
Существует дистрибутивная подрешетка $\mathcal{L}'$, содержащая $a_1, \dots,
a_n$ (это какие-то элементы), что $|L'| \le 2^{2^n}$.
\end{lmm}
\begin{proof}
\emph{(В исходном конспекте отмечено: «еще не осмыслено, просто скопировано с
теха».)}

Пусть $\mathcal{L}' = \{\varphi(a_1, \dots, a_n) \mid \varphi$ составлено из
$(+)$ и $(\cdot)\}$. Заметим, что если $p, q \in L'$, то $p \star_{\mathcal{L}} q
\in L'$ (так как $\varphi_p(\vec a) \star \varphi_q(\vec a) = \psi(\vec a)$).
Также ясно, что если $\sup_L\{p, q\} \in L'$ (или $\inf_L\{p, q\} \in L'$), то $p
\star_{\mathcal{L}} q = p \star_{\mathcal{L}'} q$. Значит, $\mathcal{L}'$ также
дистрибутивна. Построим «ДНФ»:
\[
\varphi(a_1, \dots, a_n) = \sum_{\text{Кн} \in \text{ДНФ}(\varphi)} \prod_{i \in \text{Кн}} a_i
\]
Всего не больше $2^n$ возможных компонент и $2^{2^n}$ возможных формул
$\varphi(\vec a)$.
\end{proof}

\begin{thm}
Если $\nvdash \alpha$ в ИИВ, то существует $\mathcal{G}$, что $\mathcal{G}
\not\models \alpha$, причем $|\mathcal{G}| \le 2^{2^{|\alpha|+2}}$.
\end{thm}
\begin{proof}
Используем полноту ИИВ относительно алгебр Гейтинга, если $\alpha$ где-то не
доказуема, то существует алгебра Гейтинга $\mathcal{H}$, в которой $\alpha$ не
истинна.

Рассмотрим подформулы $\alpha$. Пусть $\varphi_1, \dots, \varphi_n$~--- все
подформулы формулы $\alpha$ (включая саму $\alpha$). Число подформул $n \le
|\alpha|$. Пусть $\mathcal{G}$~--- дистрибутивная подрешетка $\mathcal{H}$,
построенная по $\llbracket\varphi_1\rrbracket, \dots, \llbracket\varphi_n\rrbracket$,
$0$ и $1$. Такая существует по лемме. $\mathcal{G}$ алгебра Гейтинга, и можно
показать, что $\mathcal{G} \not\models \alpha$. По лемме $|G| \le 2^{2^{n+2}}
\le 2^{2^{|\alpha|+2}}$.
\end{proof}

\begin{thm}
ИИВ разрешимо.
\end{thm}
\begin{proof}
Теория разрешима, если существует эффективный алгоритм (программа), который за
конечное время может определить, является ли любая заданная формула теоремой
этой теории или нет.

Давайте по формуле $\alpha$ построим все возможные алгебры Гейтинга $\mathcal{G}$
размера не больше $2^{2^{|\alpha|+37}}$, если $\mathcal{G} \models \alpha$ во
всех них, то $\proves \alpha$.
\end{proof}

\section{Категорические силлогизмы и исчисление предикатов}
\begin{dfn}
Категорический силлогизм~--- умозаключение, соединяющее три термина в двух
посылках и заключении.
\end{dfn}
\begin{dfn}[Термины]
\begin{itemize}
\item Предикат (больший термин, $P$)~--- термин, являющийся сказуемым в
  заключении.
\item Субъект (меньший термин, $S$)~--- термин, являющийся подлежащим в
  заключении.
\item Средний термин ($M$)~--- термин, присутствующий в обеих посылках, но
  отсутствующий в заключении.
\end{itemize}
\end{dfn}
\begin{dfn}[Типы высказываний (соотношений)]
\begin{itemize}
\item $A$ (общеутвердительное): Все $S$ суть $P$. Обозначение: $SaP$.
\item $I$ (частноутвердительное): Некоторые $S$ суть $P$. Обозначение: $SiP$.
\item $E$ (общеотрицательное): Ни один $S$ не есть $P$. Обозначение: $SeP$.
\item $O$ (частноотрицательное): Некоторые $S$ не суть $P$. Обозначение: $SoP$.
\end{itemize}
\end{dfn}
\begin{dfn}[Фигуры категорического силлогизма]
Расположение терминов в посылках:
\[
\begin{array}{l|c|c|c|c}
 & \text{Фигура 1} & \text{Фигура 2} & \text{Фигура 3} & \text{Фигура 4}\\\hline
\text{Большая посылка:} & M\text{---}P & P\text{---}M & M\text{---}P & P\text{---}M\\
\text{Меньшая посылка:} & S\text{---}M & S\text{---}M & M\text{---}S & M\text{---}S\\
\text{Заключение:} & S\text{---}P & S\text{---}P & S\text{---}P & S\text{---}P
\end{array}
\]
\end{dfn}
\begin{dfn}
Модус~--- комбинация типов высказываний ($A$, $I$, $E$, $O$) для большой
посылки, малой посылки и заключения в определённой фигуре (например, $AAA$ в
фигуре 1).
\end{dfn}
\begin{dfn}[Типы модусов]
\begin{itemize}
\item Сильные (правильные)~--- модусы, всегда дающие верное заключение при
  истинных посылках.
\item Слабые~--- модусы, выводящие частное заключение, когда возможно общее
  (например, вывод $SiP$ вместо $SaP$).
\item Неправильные~--- модусы, не гарантирующие истинность заключения при
  истинных посылках.
\end{itemize}
\end{dfn}
\begin{dfn}
Ограничения~--- дополнительные условия корректности для некоторых модусов
(например, требование непустоты среднего термина $M$). Или например
существования единорогов.
\end{dfn}
\begin{dfn}
Исчисление предикатов~--- формальная система для выражения высказываний с учётом
предикатов, кванторов, предметных переменных.
\end{dfn}
\begin{dfn}
Язык исчисления предикатов включает:
\begin{itemize}
\item Предметные выражения (термы):
  \begin{itemize}
  \item Предметные переменные ($x, y, \dots$).
  \item Функциональные символы ($f, g, \dots$), включая константы ($0, 1,
    \dots$).
  \end{itemize}
\item Логические выражения (формулы):
  \begin{itemize}
  \item Предикатные символы ($P, Q, \dots$), включая пропозициональные
    переменные.
  \item Логические связки: $\wedge, \vee, \to, \neg$.
  \item Кванторы: $\forall$ (всеобщность), $\exists$ (существование).
  \end{itemize}
\end{itemize}
\end{dfn}
\begin{dfn}
Метаязык~--- язык, на котором говорят об объектном языке (исчислении
предикатов), использует метапеременные для формул, термов и т.д.
\end{dfn}
\begin{dfn}
Вхождение подформулы в формулу~--- позиция первого символа этой подформулы в
формуле.
\end{dfn}
\begin{dfn}
Свободное вхождение переменной~--- вхождение переменной $x$ в формулу $\varphi$,
которое не находится в области действия квантора по $x$.
\end{dfn}
\begin{dfn}
Множество свободных переменных $FV(\varphi)$~--- множество всех переменных,
имеющих свободные вхождения в $\varphi$.
\end{dfn}
\begin{dfn}
Подстановка терма $\theta$ вместо переменной $x$ в формулу $\varphi$
($\varphi[x := \theta]$)~--- замена всех свободных вхождений $x$ в $\varphi$ на
$\theta$.
\end{dfn}
\begin{dfn}
Свобода для подстановки~--- терм $\theta$ свободен для подстановки вместо $x$ в
$\varphi$, если после подстановки никакая переменная из $\theta$ не становится
связанной в $\varphi$.
\end{dfn}
\begin{dfn}[Аксиоматика исчисления предикатов]
\begin{itemize}
\item Схемы аксиом (1)--(10) классического исчисления высказываний.
\item $(\forall x.\varphi) \to \varphi[x := \theta]$ (при условии, что $\theta$
  свободен для подстановки вместо $x$ в $\varphi$).
\item $\varphi[x := \theta] \to \exists x.\varphi$ (при том же условии).
\item Правила вывода:
  \begin{itemize}
  \item Modus Ponens (M.P.): Из $\varphi$ и $\varphi \to \psi$ выводимо $\psi$.
  \item Правило для $\forall$: Если $\proves \varphi \to \psi$, то $\proves
    \varphi \to \forall x.\psi$ (при условии, что $x$ не входит свободно в
    $\varphi$).
  \item Правило для $\exists$: Если $\proves \psi \to \varphi$, то $\proves
    \exists x.\psi \to \varphi$ (при том же условии).
  \end{itemize}
\end{itemize}
\end{dfn}
\todo{доказательства свойств силлогизмов.}

\section{Теория моделей и теоремы исчисления предикатов}
\begin{dfn}[Оценка]
Упорядоченная четвёрка $\langle D, F, P, E\rangle$, где:
\begin{itemize}
\item $D$~--- предметное множество~--- непустое множество объектов ($D \ne
  \varnothing$), из которого берутся значения для предметных переменных.
\item $F$~--- интерпретация функциональных символов: $F_{f^n} : D^n \to D$.
\item $P$~--- интерпретация предикатных символов: $P_{T^n} : D^n \to V$, где $V
  = \{\True, \False\}$.
\item $E$~--- оценка предметных переменных. $E(x) \in D$ для любой переменной
  $x$.
\end{itemize}
\end{dfn}
\begin{dfn}
Формула $\varphi$ общезначима, если она истинна при любой оценке (при любых $D$,
$F$, $P$, $E$).
\end{dfn}
\begin{dfn}
Формула $\varphi$ является логическим следствием множества формул $\Gamma$, если
$\varphi$ истинна во всех моделях, в которых истинны все формулы из $\Gamma$.
\end{dfn}

\begin{thm}[о дедукции в ИП]
Если $\Gamma \proves \alpha \to \beta$, то $\Gamma, \alpha \proves \beta$. Если
$\Gamma, \alpha \proves \beta$ и в доказательстве не применяются правила для
кванторов по свободным переменным из $\alpha$, то $\Gamma \proves \alpha \to
\beta$.
\end{thm}
\begin{proof}
($\Leftarrow$)~--- как в КИВ. ($\Rightarrow$)~--- та же схема, два новых случая.

Перестроим: $\delta_1, \delta_2, \dots, \delta_n \equiv \beta$ в $\alpha \to
\delta_1, \alpha \to \delta_2, \dots, \alpha \to \delta_n$. Дополним: обоснуем
$\alpha \to \delta_n$, если предыдущие уже обоснованы.

Два новых похожих случая: правила для $\forall$ и $\exists$. Рассмотрим
$\forall$. Доказываем $(n)\ \alpha \to \psi \to \forall x\varphi$ (правило для
$\forall$), значит, доказано $(k)\ \alpha \to \psi \to \varphi$.
\begin{align*}
&(n-0.9)\ (\alpha \to \psi \to \varphi) \to (\alpha \wedge \psi) \to \varphi &&\text{Т. о полноте КИВ}\\
&(n-0.8)\ (\alpha \wedge \psi) \to \varphi &&\text{M.P. } k, n-0.9\\
&(n-0.4)\ (\alpha \wedge \psi) \to \forall x.\varphi &&\text{Правило для } \forall,\ n-0.8\\
&(n-0.2)\ ((\alpha \wedge \psi) \to \forall x.\varphi) \to (\alpha \to \psi \to \forall x.\varphi) &&\text{Т. о полноте КИВ}\\
&(n)\ \alpha \to \psi \to \forall x.\varphi &&\text{M.P. } n-0.4,\ n-0.2
\end{align*}
\end{proof}

\begin{thm}[Корректность подстановки]
Если $\theta$ свободен для подстановки вместо $x$ в $\varphi$, то
$\llbracket\varphi\rrbracket_{x := \llbracket\theta\rrbracket} =
\llbracket\varphi[x := \theta]\rrbracket$.
\end{thm}
\begin{proof}
Будем делать индукцией по структуре $\varphi$.

\emph{База:} $\varphi$ не имеет кванторов~--- очевидно.

\emph{Индукционный переход:} пусть справедливо для $\varphi$. Покажем для
$\varphi = \forall y.\psi$.
\begin{itemize}
\item $x = y$ либо $x \notin FV(\psi)$. Тогда:
  $\llbracket\forall y.\psi\rrbracket_{x := \llbracket\theta\rrbracket} =
  \llbracket\forall y.\psi\rrbracket = \llbracket(\forall y.\psi)[x := \theta]\rrbracket$.
\item $x \ne y$. Тогда
  \[
  \llbracket\forall y.\psi\rrbracket_{x=\llbracket\theta\rrbracket}
  \overset{1}{=} \llbracket\psi\rrbracket_{y \in D,\, x := \llbracket\theta\rrbracket}
  \overset{2}{=} \llbracket\psi\rrbracket_{x := \llbracket\theta\rrbracket,\, y \in D}
  = \llbracket\psi[x := \theta]\rrbracket_{y \in D}
  = \llbracket(\forall y.\psi)[x := \theta]\rrbracket
  \]
\end{itemize}
\end{proof}

\begin{thm}[Корректность исчисления предикатов]
Если $\Gamma \proves \alpha$ и в доказательстве не используются кванторы по
свободным переменным из $FV(\Gamma)$, то $\Gamma \models \alpha$.
\end{thm}
\begin{proof}
Фиксируем $D$, $F$, $P$. Делаем индукцию по длине доказательства $\alpha$: при
любом $E$ выполнено $\Gamma \models \alpha$ при длине доказательства $n$,
покажем для $n+1$. Здесь просто аккуратно пользуемся прошлой теоремой.
\end{proof}

\section{Теорема о полноте ИП. Начало}
\begin{dfn}
$\Gamma$~--- непротиворечивое множество формул, если $\Gamma \nvdash \alpha
\amp \neg\alpha$, для любого $\alpha$.
\end{dfn}
\begin{dfn}
$\Gamma$~--- полное непротиворечивое множество замкнутых бескванторных формул,
если:
\begin{itemize}
\item $\Gamma$ содержит только замкнутые бескванторные формулы
\item если $\alpha$~--- некоторая замкнутая бескванторная формула, то либо
  $\alpha \in \Gamma$, либо $\neg\alpha \in \Gamma$
\end{itemize}
\end{dfn}
\begin{dfn}
$\Gamma$~--- полное непротиворечивое множество замкнутых формул, если:
\begin{itemize}
\item $\Gamma$ содержит только замкнутые формулы
\item если $\alpha$~--- некоторая замкнутая формула, то либо $\alpha \in
  \Gamma$, либо $\neg\alpha \in \Gamma$
\end{itemize}
\end{dfn}

\begin{lmm}
Пусть $\Gamma$ непротиворечивое множество замкнутых (бескванторных) формул.
Тогда какова бы ни была замкнутая (бескванторная) формула $\varphi$, хотя бы
$\Gamma \cup \{\varphi\}$ или $\Gamma \cup \{\neg\varphi\}$ непротиворечиво.
\end{lmm}
\begin{proof}
Из непротиворечивости $\Gamma$ следует, что непротиворечиво хотя бы одно из
$\Gamma \cup \{\varphi\}$ или $\Gamma \cup \{\neg\varphi\}$. Если бы оба были
противоречивы, то из $\Gamma$ выводилось бы как $\varphi \to (A \wedge \neg A)$,
так и $\neg\varphi \to (A \wedge \neg A)$. По законам логики из этого следовало
бы, что $\Gamma \proves A \wedge \neg A$, что противоречит непротиворечивости
$\Gamma$.
\end{proof}

\begin{thm}[о пополнении]
Пусть $\Gamma$~--- непротиворечивое множество замкнутых (бескванторных) формул.
Тогда найдётся полное непротиворечивое множество замкнутых (бескванторных)
формул $\Delta$, такое что $\Gamma \subseteq \Delta$.
\end{thm}
\begin{proof}
\begin{enumerate}
\item Занумеруем все замкнутые (бескванторные) формулы сигнатуры (их счётное
  количество): $\varphi_1, \varphi_2, \dots$
\item Построим последовательность расширяющихся множеств $\{\Gamma_i\}$
  индуктивно:
  \begin{itemize}
  \item $\Gamma_0 = \Gamma$.
  \item $\Gamma_{i+1} = \Gamma_i \cup \{\varphi_i\}$, если это множество
    непротиворечиво. В противном случае $\Gamma_{i+1} = \Gamma_i \cup
    \{\neg\varphi_i\}$ (это верно из леммы).
  \end{itemize}
\item Положим $\Delta = \bigcup_i \Gamma_i$.
\item Докажем, что $\Delta$ непротиворечиво. Предположим противное: $\Delta
  \proves \alpha \wedge \neg\alpha$. Доказательство имеет конечную длину и
  использует конечное число гипотез $\{\delta_1, \dots, \delta_n\} \subseteq
  \Delta$. Пусть каждая $\delta_i$ принадлежит $\Gamma_{d_i}$. Тогда все они
  содержатся в $\Gamma_m$, где $m = \max(d_1, \dots, d_n)$. Следовательно,
  $\Gamma_m \proves \alpha \wedge \neg\alpha$, что противоречит построению
  $\Gamma_m$ как непротиворечивого множества.
\end{enumerate}
\end{proof}

\begin{dfn}
Моделью для множества формул $F$ назовем такую модель $\mathcal{M}$, что при
всяком $\varphi \in F$ выполнено $\llbracket\varphi\rrbracket_{\mathcal{M}} =
\True$.
\end{dfn}
\begin{thm}
Любое непротиворечивое множество замкнутых формул имеет модель.
\end{thm}
\begin{proof}
Построим. Пусть $M$~--- полное непротиворечивое множество замкнутых
бескванторных формул. Тогда модель $\mathcal{M}$ задается так:
\begin{enumerate}
\item $D$~--- множество всевозможных предметных выражений без предметных
  переменных.
\item $\llbracket f(\theta_1, \dots, \theta_n)\rrbracket = \text{"f("} +
  \llbracket\theta_1\rrbracket + \text{", "} + \dots + \text{", "} +
  \llbracket\theta_n\rrbracket + \text{")"}$
\item $\llbracket P(\theta_1, \dots, \theta_n)\rrbracket = \True$, если
  $P(\theta_1, \dots, \theta_n) \in M$
\item Так как $D \ne \varnothing$, то найдется $z \in D$. Тогда
  $\llbracket x\rrbracket = z$. Это ничему не помешает, так как формулы замкнуты.
\end{enumerate}
\begin{lmm}
$\mathcal{M} \models \varphi$ тогда и только тогда, когда $\varphi \in M$.
\end{lmm}
\emph{Доказательство:} Индукция по длине формулы $\varphi$.

\emph{База:} предикат.

\emph{Индукционный переход:} Нам надо показать, что $\mathcal{M} \models \alpha
\star \beta \Leftrightarrow \alpha \star \beta \in M$. Это делается просто
перебором оценок $\alpha$ и $\beta$.

\emph{Концовка.} Пусть $M$~--- непротиворечивое множество замкнутых
бескванторных формул. По теореме о пополнении существует $M'$~--- полное
непротиворечивое множество замкнутых бескванторных формул, что $M \subseteq
M'$. По лемме $M'$ имеет модель, эта модель подойдёт для $M$.
\end{proof}

\section{Теорема о полноте ИП, продолжение}
\begin{dfn}
Формула $\varphi$ имеет поверхностные кванторы (находится в предваренной форме),
если
\[ \varphi := \forall x.\varphi \mid \exists x.\varphi \mid \tau \]
где $\tau$~--- формула без кванторов. То есть: $\forall x\exists y\exists
z\forall t : \tau$.
\end{dfn}
\begin{thm}
Для любой замкнутой формулы $\varphi$ найдется такая формула $\psi$ с
поверхностными кванторами, что $\proves \psi \to \varphi$ и $\proves \varphi \to
\psi$.
\end{thm}
\begin{dfn}
Сколемизация: \todo{?????}
\end{dfn}
\begin{thm}[Теорема Гёделя о полноте исчисления предикатов]
Если $M$ замкнутое непротиворечивое множество формул, то оно имеет модель.
\end{thm}
\begin{proof}
Построим по $M$ множество $M'$ с поверхностными кванторами. По $M'$ построим
непротиворечивое множество замкнутых бескванторных, дополним его до полного,
построим для него модель. Тадам.
\end{proof}
\begin{thm}[Следствие (теорема о полноте)]
Исчисление предикатов полно.
\end{thm}
\begin{proof}
Доказательство от противного.
\end{proof}

\section{Машина Тьюринга}
\begin{dfn}
Машина Тьюринга~--- система, задаваемая:
\begin{itemize}
\item Внешним алфавитом: конечный набор символов $q_1, \dots, q_n$, включая
  выделенный символ-заполнитель $q_\varepsilon$.
\item Внутренний алфавит (состояний): конечный набор состояний $s_1, \dots,
  s_k$, включая:
  \begin{itemize}
  \item $s_s$~--- начальное состояние,
  \item $s_f$~--- допускающее состояние,
  \item $s_r$~--- отвергающее состояние.
  \end{itemize}
\item Таблица переходов: функция, которая для каждой пары (текущее состояние
  $s$, текущий символ $q$) задаёт тройку (новое состояние $s'$, символ для
  записи $q'$, направление движения головки $\{\leftarrow, \rightarrow,
  \cdot\}$).
\end{itemize}
\end{dfn}
\begin{dfn}
Язык~--- произвольное множество конечных слов (строк) над некоторым конечным
алфавитом.
\end{dfn}
\begin{dfn}[Разрешимый язык]
Язык $L$ называется разрешимым, если существует машина Тьюринга, которая для
любого входного слова $w$:
\begin{itemize}
\item Останавливается в допускающем состоянии ($s_f$), если $w \in L$.
\item Останавливается в отвергающем состоянии ($s_r$), если $w \notin L$.
\item Такая машина называется разрешающей для языка $L$.
\end{itemize}
\end{dfn}
\begin{dfn}[Задача об останове (проблема останова)]
Задача определения по описанию машины Тьюринга $M$ и входному слову $w$,
остановится ли машина $M$ на входе $w$.
\end{dfn}

\begin{thm}[о неразрешимости задачи об останове]
Язык останавливающихся машин Тьюринга неразрешим.
\end{thm}
\begin{proof}
Схема доказательства (от противного): Предположим, существует разрешающая машина
$S$, которая определяет, остановится ли машина $x$ на входе $y$. Построим машину
$W$, которая на входе $x$ делает следующее:
\begin{enumerate}
\item Запускает $S$ на паре $(x, x)$ (т.е.\ спрашивает, остановится ли $x$ на
  своём собственном коде).
\item Если $S$ отвечает «да» (остановится), то $W$ уходит в бесконечный цикл.
\item Если $S$ отвечает «нет» (не остановится), то $W$ останавливается и
  возвращает 1.
\end{enumerate}
Если теперь подать код $W$ на вход самой $W$, то получим противоречие: $S(W,
W)$ должна быть истинна тогда и только тогда, когда она ложна.
\end{proof}

\begin{thm}
Язык всех доказуемых формул исчисления предикатов неразрешим.
\end{thm}
\begin{proof}
Пусть существует машина Тьюринга, разрешающая любую формулу. На ее основе тогда
несложно построить некоторую машину Тьюринга, перестраивающую любую машину $S$ (с
допускающим состоянием $s_f$ и входом $y$) в ее ограничения $C$ и разрешающую
формулу ИП $C \to \exists w_l.\exists w_r.F_{S y}(w_l, w_r, s_f)$. Эта машина
разрешит задачу останова, а такого не может быть. \todo{}
\end{proof}

\section{Аксиоматика Пеано, арифметизация}
\begin{dfn}[Представление чисел через натуральные]
\begin{itemize}
\item Целые $Z = \{\langle x, y\rangle \mid x, y \in \mathbb{N}_0\}$.
  Интуиция: $\langle x, y\rangle = x - y$; $\langle a, b\rangle + \langle c,
  d\rangle = \langle a + c, b + d\rangle$; $\langle a, b\rangle - \langle c,
  d\rangle = \langle a + d, b + c\rangle$. Пусть $\langle a, b\rangle \equiv
  \langle c, d\rangle$, если $a + d = b + c$. Тогда $\mathbb{Z} = Z/{\equiv}$;
  $0 = [\langle 0, 0\rangle]$, $1 = [\langle 1, 0\rangle]$, $-7 = [\langle 0,
  7\rangle]$.
\item Рациональные $Q = \mathbb{Z} \times \mathbb{N}$~--- множество всех простых
  дробей. $\langle p, q\rangle$~--- то же, что $p/q$; $\langle p_1, q_1\rangle
  \equiv \langle p_2, q_2\rangle$, если $p_1 q_2 = p_2 q_1$. $\mathbb{Q} =
  Q/{\equiv}$.
\item Вещественные $X = \{A, B\}$, где $A, B \subseteq \mathbb{Q}$~---
  дедекиндово сечение, если: (a) $A \cup B = \mathbb{Q}$; (b) если $a \in A$,
  $x \in \mathbb{Q}$ и $x \le a$, то $x \in A$; (c) если $b \in B$, $x \in
  \mathbb{Q}$ и $b \le x$, то $x \in B$; (d) $A$ не содержит наибольшего.
  $\mathbb{R}$~--- множество всех возможных дедекиндовых сечений.
  $\sqrt{2} = \{\{x \in \mathbb{Q} \mid x < 0 \vee x^2 < 2\}, \{x \in
  \mathbb{Q} \mid x > 0 \amp x^2 > 2\}\}$.
\end{itemize}
\end{dfn}
\begin{dfn}[Аксиоматика Пеано]
$N$ (более точно, $\langle N, 0, (')\rangle$) соответствует аксиоматике Пеано,
если следующее определено/выполнено:
\begin{enumerate}
\item Операция «штрих» $(') : N \to N$, причём нет $a, b \in N$, что $a \ne b$,
  но $a' = b'$. Если $x = y'$, то $x$ назовём следующим за $y$, а $y$~---
  предшествующим $x$.
\item Константа $0 \in N$: нет $x \in N$, что $x' = 0$.
\item Индукция. Каково бы ни было свойство («предикат») $P : N \to V$, если:
  (a) $P(0)$; (b) при любом $x \in N$ из $P(x)$ следует $P(x')$, то при любом
  $x \in N$ выполнено $P(x)$.
\end{enumerate}
Как построить? Например, в стиле алгебры Линденбаума:
\begin{enumerate}
\item $N$~--- язык, порождённый грамматикой $\nu ::= 0 \mid \nu\text{«'»}$
\item $0$~--- это «0», $x'$~--- это $x + \text{«'»}$
\end{enumerate}
\end{dfn}
\begin{dfn}[Арифметические операции (возведение в степень) в аксиоматике Пеано]
\begin{itemize}
\item цифры: $1 = 0'$, $2 = 0''$, $3 = 0'''$
\item сложение: $a + b = \begin{cases} a, & b = 0\\ (a + c)', & b = c'\end{cases}$
\item умножение: $a \cdot b = \begin{cases} 0, & b = 0\\ a \cdot c + a, & b = c'\end{cases}$
\item возведение в степень: $a^b = \dots$
\end{itemize}
\end{dfn}
\begin{dfn}
Доказательство коммутативности сложения.
\end{dfn}
\begin{dfn}
Порядок теории (0, 1, 2).
\end{dfn}
\begin{dfn}
Теорией первого порядка назовём исчисление предикатов с дополнительными
(«нелогическими» или «математическими»):
\begin{itemize}
\item предикатными и функциональными символами;
\item аксиомами.
\end{itemize}
Сущности, взятые из исходного исчисления предикатов, назовём логическими.
\end{dfn}
\begin{dfn}
Формальная арифметика~--- теория первого порядка, со следующими добавленными
нелогическими\dots
\begin{itemize}
\item двухместными функциональными символами $(+)$, $(\cdot)$; одноместным
  функциональным символом $(')$, нульместным функциональным символом $0$;
\item двухместным предикатным символом $(=)$;
\item восемью нелогическими аксиомами:
  \begin{align*}
  &(A1)\ a = b \to a = c \to b = c\\
  &(A2)\ a = b \to a' = b'\\
  &(A3)\ a' = b' \to a = b\\
  &(A4)\ \neg a' = 0\\
  &(A5)\ a + 0 = a\\
  &(A6)\ a + b' = (a + b)'\\
  &(A7)\ a \cdot 0 = 0\\
  &(A8)\ a \cdot b' = a \cdot b + a
  \end{align*}
\item нелогической схемой аксиом индукции $\psi[x := 0] \amp (\forall x.\psi \to
  \psi[x := x']) \to \psi$ с метапеременными $x$ и $\psi$.
\end{itemize}
\end{dfn}
\begin{dfn}[Доказательство $a = a$]
Пусть $\top ::= 0 = 0 \to 0 = 0 \to 0 = 0$, тогда:
\begin{align*}
&(1)\ a = b \to a = c \to b = c &&\text{(Акс.\ А1)}\\
&(2)\ (a = b \to a = c \to b = c) \to \top \to (a = b \to a = c \to b = c) &&\text{(Сх. акс. 1)}\\
&(3)\ \top \to (a = b \to a = c \to b = c) &&\text{(M.P. 1, 2)}\\
&(4)\ \top \to (\forall c.a = b \to a = c \to b = c) &&\text{(Введ. $\forall$)}\\
&(5)\ \top \to (\forall b.\forall c.a = b \to a = c \to b = c) &&\text{(Введ. $\forall$)}\\
&(6)\ \top \to (\forall a.\forall b.\forall c.a = b \to a = c \to b = c) &&\text{(Введ. $\forall$)}\\
&(7)\ \top &&\text{(Сх. акс. 1)}\\
&(8)\ (\forall a.\forall b.\forall c.a = b \to a = c \to b = c) &&\text{(M.P. 7, 6)}\\
&(9)\ (\forall a.\forall b.\forall c.a = b \to a = c \to b = c) \to (\forall b.\forall c.a + 0 = b \to a + 0 = c \to b = c) &&\text{(Сх. акс. 11)}\\
&(10)\ \forall b.\forall c.a + 0 = b \to a + 0 = c \to b = c &&\text{(M.P. 8, 9)}\\
&(12)\ \forall c.a + 0 = a \to a + 0 = c \to a = c &&\text{(M.P. 10, 11)}\\
&(14)\ a + 0 = a \to a + 0 = a \to a = a &&\text{(M.P. 12, 13)}\\
&(15)\ a + 0 = a &&\text{(Акс. А5)}\\
&(16)\ a + 0 = a \to a = a &&\text{(M.P. 15, 14)}\\
&(17)\ a = a &&\text{(M.P. 15, 16)}
\end{align*}
\end{dfn}
\begin{dfn}[Арифметизация математики, формализация категорических силлогизмов, предложенная Лейбницем]
\begin{itemize}
\item Любой термин~--- пара взаимно простых чисел $+a - b$. Например, мудрый~---
  $+70 - 33$, благочестивый~--- $+10 - 3$.
\item Общеутвердительное предложение (каждый $+a - b$ есть $+c - d$): $a : c$ и
  $b : d$. Всякий мудрый есть благочестивый ($70 = 10 \cdot 7$, $33 = 3 \cdot
  11$).
\item Частноотрицательное предложение~--- не верно общеутвердительное.
\item Общеотрицательное предложение~--- когда $a, d$ или $b, c$ имеют общий
  делитель, отличный от 1: Ни один благочестивый ($+10 - 3$) не есть несчастный
  ($+5 - 14$), так как $10 = 2 \cdot 5$ и $14 = 2 \cdot 7$.
\end{itemize}
\end{dfn}

\section{Функции}
\begin{dfn}[Базовые примитивы (исходные функции)]
\begin{itemize}
\item Ноль ($Z$): $Z(x) = 0$.
\item Инкремент ($N$): $N(x) = x + 1$.
\item Проекция ($U$): $U^n_k(x_1, \dots, x_n) = x_k$ (возвращает $k$-й аргумент).
\item Подстановка ($S$): Если $g$~--- функция от $k$ аргументов, а $f_1, \dots,
  f_k$~--- функции от $n$ аргументов, то $\Sh{g, f_1, \dots, f_k}(\vec x) =
  g(f_1(\vec x), \dots, f_k(\vec x))$.
\end{itemize}
\end{dfn}
\begin{dfn}[Примитивная рекурсия]
Если $f$~--- функция от $n$ аргументов, а $g$~--- функция от $n+2$ аргументов, то
функция $h = \Rr{f, g}$ от $n+1$ аргумента определяется так:
\begin{itemize}
\item $h(\vec x, 0) = f(\vec x)$
\item $h(\vec x, y+1) = g(\vec x, y, h(\vec x, y))$
\end{itemize}
\end{dfn}
\begin{dfn}
Примитивно-рекурсивная функция~--- функция, которую можно получить из базовых
примитивов ($Z$, $N$, $U$) с помощью конечного числа применений операций
подстановки ($S$) и примитивной рекурсии ($R$).
\end{dfn}
\begin{dfn}[Операция минимизации ($M$)]
Если $f$~--- функция от $n+1$ аргумента, то $\Mn{f}(\vec x) = \min\{y \mid
f(\vec x, y) = 0\}$, при условии, что такое $y$ существует. Если для всех $y$
значение $f(\vec x, y) > 0$, то $\Mn{f}(\vec x)$ не определено.
\end{dfn}
\begin{dfn}
Общерекурсивная функция~--- функция, которую можно получить из базовых
примитивов с помощью конечного числа применений операций подстановки ($S$),
примитивной рекурсии ($R$) и минимизации ($M$).
\end{dfn}

\begin{lmm}
Функция Аккермана не выражена в примитивно рекурсивных.
\[
A(m, n) =
\begin{cases}
n + 1, & m = 0\\
A(m-1, 1), & m > 0,\ n = 0\\
A(m-1, A(m, n-1)), & m > 0,\ n > 0
\end{cases}
\]
\end{lmm}
\begin{proof}
Для этого докажем немного другую теорему.
\end{proof}
\begin{thm}
Пусть $f(\vec x)$~--- примитивно рекурсивная. Тогда найдется $k$, что $f(\vec x)
< A(k, \max(\vec x))$.
\end{thm}
\begin{proof}
Индукция по структуре $f$.
\end{proof}

\begin{dfn}
Будем говорить, что отношение $R \subseteq \mathbb{N}_0^n$ выразимо в ФА, если
существует формула $\rho$, что:
\begin{itemize}
\item если $\langle a_1, \dots, a_n\rangle \in R$, то $\proves \rho(a_1, \dots,
  a_n)$
\item если $\langle a_1, \dots, a_n\rangle \notin R$, то $\proves \neg\rho(a_1,
  \dots, a_n)$
\end{itemize}
\end{dfn}
\begin{thm}
Отношение «равно» выразимо в Ф.А.: $r = \{\langle x, x\rangle \mid x \in
\mathbb{N}_0\}$.
\end{thm}
\begin{dfn}
Будем говорить, что функция $f : \mathbb{N}_0^n \to \mathbb{N}_0$ представима в
ФА, если существует формула $\varphi$, что:
\begin{itemize}
\item если $f(a_1, \dots, a_n) = u$, то $\proves \varphi(a_1, \dots, a_n, u)$
\item если $f(a_1, \dots, a_n) \ne u$, то $\proves \neg\varphi(a_1, \dots, a_n,
  u)$
\item для всех $a_i \in \mathbb{N}_0$ выполнено $\proves (\exists
  x.\varphi(a_1, \dots, a_n, x)) \amp (\forall p, \forall q.\ \varphi(a_1,
  \dots, a_n, p) \amp \varphi(a_1, \dots, a_n, q) \to p = q)$
\end{itemize}
\end{dfn}
\begin{dfn}
Для отношения $R$ его характеристическая функция $\chi_R$ возвращает 1, если
аргумент принадлежит $R$, и 0 иначе. Отношение выразимо тогда и только тогда,
когда представима его характеристическая функция.
\end{dfn}

\section{Представимость примитивов, бета-функции}
\begin{dfn}[Представимость базовых примитивов]
\begin{itemize}
\item Функция ноль ($Z$): $Z(x) = 0$. Представима формулой $\zeta(x, y) := (x =
  x) \wedge (y = 0)$.
\item Функция следования ($N$): $N(x) = x + 1$. Представима формулой $\nu(x, y)
  := y = x'$.
\item Функция проекции ($U^k_n$): $U^k_n(x_1, \dots, x_n) = x_k$. Представима
  формулой $\upsilon(x_1, \dots, x_n, y) := y = x_k$.
\item Операция подстановки: Если функции $g(y_1, \dots, y_m)$ и $h_1(\vec x),
  \dots, h_m(\vec x)$ представимы формулами $\gamma$ и $\eta_1, \dots, \eta_m$
  соответственно, то функция $f(\vec x) = g(h_1(\vec x), \dots, h_m(\vec x))$
  представима формулой:
  \[
  \varphi(\vec x, y) := \exists z_1 \dots \exists z_m\,(\gamma(z_1, \dots, z_m, y) \wedge \eta_1(\vec x, z_1) \wedge \dots \wedge \eta_m(\vec x, z_m))
  \]
\end{itemize}
\end{dfn}
\begin{dfn}[$\beta$-функция Гёделя]
$\beta(b, c, i)$ = остаток от деления $b$ на $(1 + c \cdot (i + 1))$.
\end{dfn}
\begin{thm}
Если $a_0, \dots, a_n \in \mathbb{N}_0$, то найдутся такие $b, c \in
\mathbb{N}_0$, что $a_i = \beta(b, c, i)$.
\end{thm}
\begin{thm}
Примитив $\Rr{f, g}$ представим в Ф.А.\ формулой $\rho(x_1, \dots, x_n, y, a)$:
\begin{multline*}
\exists b.\exists c.\,(\exists a_0.\,\beta(b, c, 0, a_0) \amp \phi(x_1, \dots x_n, a_0))\\
\amp \forall k.\,k < y \to \exists d.\exists e.\,\beta(b, c, k, d) \amp \beta(b, c, k', e) \amp \gamma(x_1, \dots x_n, k, d, e)\\
\amp \beta(b, c, y, a)
\end{multline*}
\end{thm}
\begin{thm}
Пусть функция $f : \mathbb{N}_0^{n+1} \to \mathbb{N}_0$ представима в Ф.А.\
формулой $\phi(x_1, \dots, x_n, y, r)$. Тогда примитив $\Mn{f}$ представим в
Ф.А.\ формулой:
\[
\mu(x_1, \dots, x_n, y) := \phi(x_1, \dots, x_n, y, 0) \amp \forall u.\,u < y \to \neg\phi(x_1, \dots, x_n, u, 0)
\]
\end{thm}
\begin{thm}
Если $f$~--- рекурсивная функция, то она представима в Ф.А. Индукция по
структуре $f$.
\end{thm}

\section{Гёделева нумерация. Начало пиздеца}
\begin{dfn}[Рекурсивность представимых в Ф.А. функций]
Фиксируем $f$ и $x_1, x_2, \dots, x_n$. Обозначим $y = f(x_1, x_2, \dots, x_n)$.
По представимости нам известна $\phi$, что $\proves \phi(x_1, x_2, \dots, x_n,
y)$. Давайте просто переберём все результаты и доказательства!
\begin{enumerate}
\item Закодируем доказательства натуральными числами.
\item Напишем рекурсивную функцию, проверяющую доказательства на корректность.
\item Параллельный перебор значений и доказательств: $s = 2^y \cdot 3^p$.
\end{enumerate}
Переберём все $s$, по $s$ получим $y$ и $p$. Проверим, что $p$~--- код
доказательства $\proves \phi(x_1, x_2, \dots, x_n, y)$.
\end{dfn}
\begin{thm}
Следующая функция рекурсивна:
\[
\mathrm{proof}(f, x_1, x_2, \dots, x_n, y, p) =
\begin{cases}
1, & \text{если } p \text{ --- гёделев номер доказательства}\\
0, & \text{иначе}
\end{cases}
\]
\end{thm}
\begin{dfn}
Пусть $\xi$~--- формула с единственной свободной переменной $x_1$. Тогда
$\langle\gn{\xi}, p\rangle \in W_1$, если $\proves \xi(\gn{\xi})$ и $p$~--- номер
доказательства.
\end{dfn}
\begin{dfn}
Отношение $W_1$ рекурсивно, поэтому выражено в Ф.А.\ формулой $\omega_1$ со
свободными переменными $x_1$ и $x_2$, причём:
\begin{enumerate}
\item $\proves \omega_1(\gn{\phi}, p)$, если $p$~--- гёделев номер доказательства
  самоприменения $\phi$;
\item $\proves \neg\omega_1(\gn{\phi}, p)$ иначе.
\end{enumerate}
\end{dfn}
\todo{$W_2$, $w_2$ --- что за ферзь?}

\section{Теоремы Гёделя и пиздец больного ублюдка}
\begin{dfn}
Теория называется $\omega$-непротиворечивой, если для любой формулы $\varphi(x)$
из того, что доказуемы $\varphi(0), \varphi(1), \varphi(2), \dots$ (т.е.\
$\varphi(n)$ для всех натуральных $n$), следует, что не доказуемо $\exists
x.\neg\varphi(x)$.
\end{dfn}
\begin{dfn}
Определим формулу $\sigma(x_1) := \forall p.\neg\omega_1(x_1, p)$.
\end{dfn}

\begin{thm}[Первая теорема Гёделя о неполноте арифметики]
\begin{itemize}
\item Если формальная арифметика непротиворечива, $\nvdash \sigma(\gn{\sigma})$.
\item Если формальная арифметика $\omega$-непротиворечива, $\nvdash
  \neg\sigma(\gn{\sigma})$.
\end{itemize}
\end{thm}
\begin{proof}
Напомним $\sigma(x_1) := \forall p.\neg w_1(x_1, p)$.
\begin{enumerate}
\item Пусть $\proves \sigma(\gn{\sigma})$. Тогда есть $p$~--- номер
  доказательства, тогда $\proves w_1(\gn{\sigma}, p)$. Тогда $\proves \exists
  p.w_1(\gn{\sigma}, p)$, тогда $\proves \neg\sigma(\gn{\sigma})$. Противоречие.
\item Пусть $\proves \neg\sigma(\gn{\sigma})$ \emph{(в исходном конспекте
  опечатка: написано $\proves \sigma(\gn{\sigma})$)}. То есть $\exists p.w_1(\gn{\sigma}, p)$.
  \begin{itemize}
  \item Но найдется ли натуральное число $p$, что $\proves w_1(\gn{\sigma}, p)$?
    Пусть нет. Тогда по $\omega$-непротиворечивости $\nvdash \exists
    p.\neg\neg w_1(\gn{\sigma}, p)$. То есть $p$ доказательство самоприменения:
    $\proves \omega(\gn{\sigma})$. Противоречие.
  \end{itemize}
  \todo{что что что?}
\end{enumerate}
\end{proof}

\begin{dfn}
Теория называется синтаксически полной, если для любой замкнутой формулы
$\alpha$ этой теории верно либо $\proves \alpha$, либо $\proves \neg\alpha$.
\end{dfn}
\begin{dfn}
Теория называется семантически полной (относительно своей стандартной модели),
если любая формула, истинная в стандартной модели, доказуема в теории.
\end{dfn}
\begin{thm}
Формальная арифметика (при условии её непротиворечивости) не является ни
синтаксически, ни семантически полной. Существуют формулы, которые неразрешимы
(не доказуемы и не опровержимы), и существуют истинные в стандартной модели
формулы, которые недоказуемы.
\end{thm}
\begin{proof}
Рассмотрим Ф.А.\ с классической моделью. Из теоремы Гёделя имеем $\nvdash
\sigma(\gn{\sigma})$. Рассмотрим $\sigma(\gn{\sigma}) \equiv \forall
p.\neg\omega_1(\gn{\sigma}, p)$: нет числа $p$, что $p$~--- номер доказательства
$\sigma(\gn{\sigma})$. То есть, $\llbracket\forall p.\neg w_1(\gn{\sigma},
p)\rrbracket$. То есть $\models \sigma(\gn{\sigma})$.
\end{proof}

\begin{dfn}
Пусть $\langle\gn{\xi}, p\rangle \in W_2$, если $\proves \neg\xi(\gn{\xi})$ и $p$
номер доказательства. Пусть $w_2$ выражает $W_2$ в формальной арифметике.
\end{dfn}
\begin{thm}[Первая теорема Гёделя в форме Россера]
Рассмотрим $\rho(x) = \forall p.\,w_1(x_1, p) \to \exists q.\,q \le p \amp
w_2(x_2, q)$. Тогда $\nvdash \rho(\gn{\rho})$ и $\nvdash \neg\rho(\gn{\rho})$.
Неформальный смысл: «Меня легче опровергнуть, чем доказать».
\end{thm}
\begin{dfn}
Теория первого порядка, использующая нелогические функциональные символы $0$,
$(+)$ и $(\cdot)$, нелогический предикатный символ $(=)$ и следующие
нелогические аксиомы, называется системой Робинсона.
\begin{align*}
&a = a & &a = b \to b = a\\
&a = b \to b = c \to a = c & &a = b \to a' = b'\\
&a' = b' \to a = b & &\neg 0 = a'\\
&a = b \to a + c = b + c \amp c + a = c + b & &a = b \to a \cdot c = b \cdot c \amp c \cdot a = c \cdot b\\
&\neg a = 0 \to \exists b.\,a = b' & &a + 0 = a\\
&a + b' = (a + b)' & &a \cdot 0 = 0\\
&a \cdot b' = a \cdot b + a
\end{align*}
Система Робинсона неполна: аксиомы~--- в точности утверждения, необходимые для
доказательства теорем Гёделя. Система Робинсона не имеет схем аксиом.
\end{dfn}
\begin{dfn}
Теория первого порядка, использующая нелогические функциональные символы $0$,
$1$, $(+)$, нелогический предикатный символ $(=)$ и следующие нелогические
аксиомы, называется арифметикой Пресбургера.
\begin{align*}
&\neg(0 = x + 1)\\
&x + 1 = y + 1 \to x = y\\
&x + 0 = x\\
&x + (y + 1) = (x + y) + 1\\
&(\varphi(0) \amp \forall x.\,\varphi(x) \to \varphi(x + 1)) \to \forall y.\,\varphi(y)
\end{align*}
\end{dfn}
\begin{thm}
Арифметика Пресбургера разрешима и синтаксически и семантически полна.
\end{thm}

\section{Продолжение теорем Гёделя}
\begin{dfn}
Обозначим за $\varphi(x, p)$, выражающую в формальной арифметике рекурсивное
отношение $\mathrm{Proof}$. Обозначим $\pi(x) = \exists p.\varphi(x, p)$.
(Идейно, формула доказуема.)
\end{dfn}
\begin{dfn}
$\mathrm{Consis} := \neg\pi(\gn{1 = 0})$. Неформальный смысл: формальная
арифметика непротиворечива.
\end{dfn}
\begin{thm}[Вторая теорема Гёделя о неполноте арифметики]
Если $\mathrm{Consis}$ доказуем, то формальная арифметика противоречива.
\end{thm}
\begin{proof}
Если $\mathrm{Consis}$ доказуем, то если мы покажем, что $\proves \mathrm{Consis}
\to \sigma(\gn{\sigma})$, то из первой теоремы мы проиграем.

Рассмотрим особый $\mathrm{Consis}'$:
\[ \pi'(x) = \exists p.\,\varphi(x, p) \amp \neg\varphi(\gn{1 = 0}, p) \]
\[ \mathrm{Consis}' = \neg\pi'(\gn{1 = 0}) \]
Если ФА непротиворечива, то $\llbracket\pi'(x) = \pi(x)\rrbracket$.
\todo{} \todo{} \todo{} \todo{} \todo{} \todo{}
\end{proof}
\begin{dfn}[Условия выводимости Гильберта--Бернайса--Лёба]
Будем говорить, что формула $\psi$, выражающая отношение $\mathrm{Proof}$,
формула $\pi$ и формула $\mathrm{Consis}$ соответствуют условиям
Гильберта--Бернайса--Лёба, если следующие условия выполнены для любой формулы
$\alpha$:
\begin{enumerate}
\item $\proves \alpha$ влечет $\proves \pi(\gn{\alpha})$
\item $\proves \pi(\gn{\alpha}) \to \pi(\gn{\pi(\gn{\alpha})})$
\item $\proves \pi(\gn{\alpha \to \beta}) \to \pi(\gn{\alpha}) \to \pi(\gn{\beta})$
\end{enumerate}
\end{dfn}
\begin{dfn}[Лемма об автоссылках]
Для любой формулы $\varphi(x_1)$ можно построить такую замкнутую формулу
$\alpha$ (не использующую неаксиоматических предикатных и функциональных
символов), что $\proves \varphi(\gn{\alpha}) \leftrightarrow \alpha$.
\end{dfn}
\begin{thm}[Непротиворечивость ФА]
Существует такая замкнутая формула $\gamma$, что если Ф.А.\ непротиворечива, то
$\nvdash \gamma$, а если Ф.А.\ $\omega$-непротиворечива, то и $\nvdash
\neg\gamma$ \emph{(в исходном конспекте: «и $\proves \neg\gamma$»)}.
\end{thm}
\begin{proof}
Рассмотрим $\Phi(x_1) = \neg\pi(x_1)$. Тогда по лемме об автоссылках существует
$\gamma$, что $\proves \gamma \leftrightarrow \neg\pi(\gn{\gamma})$.
\begin{itemize}
\item Предположим, что $\proves \gamma$. Тогда по вышесказанному $\nvdash
  \gamma$.
\item Предположим, что $\proves \neg\gamma$. Тогда $\proves \pi(\gn{\gamma})$,
  то есть $\proves \exists p.\psi(\gn{\gamma}, p)$. Тогда по
  $\omega$-непротиворечивости найдется $p$, что $\proves \varphi(\gn{\gamma},
  p)$, то есть $\proves \gamma$.
\end{itemize}
\end{proof}
\begin{dfn}
Теория $\mathcal{S}$~--- расширение теории $\mathcal{T}$, если из $\proves_{\mathcal{T}}
\alpha$ следует $\proves_{\mathcal{S}} \alpha$.
\end{dfn}
\begin{dfn}
Теория $\mathcal{S}$~--- рекурсивно-аксиоматизируемая, если найдется
$\mathcal{S}'$ с тем же языком, что
\begin{itemize}
\item $\proves_S \alpha$ тогда и только тогда когда $\proves_{S'} \alpha$
\item Множество аксиом $S'$ рекурсивно (????)
\end{itemize}
\end{dfn}
\begin{dfn}
$Th_S = \{\gn{\alpha} \mid \proves_S \alpha\}$, $Tr_S = \{\gn{\alpha} \mid
\llbracket\alpha\rrbracket_S = \True\}$, то есть множество доказуемых и
общезначимых в теории.
\end{dfn}
\todo{тут скипнуто 2 леммы, нужны ли они}

\begin{thm}[о неразрешимости формальной арифметики]
Если формальная арифметика непротиворечива, то формальная арифметика
неразрешима.
\end{thm}
\begin{proof}
Пусть формальная арифметика разрешима, откуда есть рекурсивная функция $f(x)$:
$f(x) = 1$ тогда и только тогда, когда $x \in Th_{\text{Ф.А}}$. То есть
$Th_{\text{Ф.А}}$ выразимо в формальной арифметике.
\end{proof}

\begin{thm}[Тарского о невыразимости истины]
Не существует $\varphi(x)$, что $\llbracket\varphi(x)\rrbracket = \True$ (в
стандартной интерпретации) тогда и только тогда, когда $x \in Tr_{\text{Ф.А}}$.
\end{thm}
\begin{proof}
Пусть теория $S$~--- формальная арифметика + аксиомы: все истинные в
стандартной интерпретации формулы. Очевидно, что $Th_S = Tr_S = Tr_{\text{Ф.А}}$.
То есть $Tr_{\text{Ф.А}}$ невыразимо в $S$.

Пусть $\varphi$ таково, что $\llbracket\varphi(x)\rrbracket = \True$ при $x \in
Tr$. Тогда $\proves \varphi(x)$, если $x \in Tr$ и $\proves \neg\varphi(x)$,
если $x \notin Tr$, тогда $Tr$ выразимо в $S$~--- противоречие. \todo{}
\end{proof}

\section{Множества}
\begin{dfn}[Определения равенства]
\begin{itemize}
\item Два множества равны, если они содержат одни и те же элементы (аксиома
  экстенсиональности).
\item Формально: $A = B \iff \forall x(x \in A \leftrightarrow x \in B)$.
\end{itemize}
\end{dfn}
\begin{dfn}[Аксиоматика Цермело--Френкеля (ZF)]
\begin{itemize}
\item Равенство «по Лейбницу»: объекты равны, если неразличимы. Если нечто ходит
  как утка, выглядит как утка и крякает как утка, то это утка.
\item Принцип объёмности: объекты равны, если состоят из одинаковых частей.
  \[ A \subseteq B \equiv \forall x.\,x \in A \to x \in B \qquad A = B \equiv A \subseteq B \amp B \subseteq A \]
\item Аксиома равенства: равные множества содержатся в одних и тех же
  множествах.
  \[ \forall x.\forall y.\forall z.\,x = y \amp x \in z \to y \in z \]
\item Аксиома пустого. Существует пустое множество $\varnothing$.
  $\exists s.\forall t.\,\neg t \in s$
\item Аксиома пары. Существует $\{a, b\}$. Каковы бы ни были два множества $a$ и
  $b$, существует множество, состоящее в точности из них.
  \[ \forall a.\forall b.\exists s.\,a \in s \amp b \in s \amp \forall c.\,c \in s \to c = a \vee c = b \]
\item Аксиома объединения: существует $\bigcup x$. Для любого непустого
  множества $x$ найдется такое множество, состоящее в точности из тех
  элементов, из которых состоят элементы $x$.
  \[ \forall x.\,(\exists y.\,y \in x) \to \exists p.\forall y.\,y \in p \leftrightarrow \exists s.\,y \in s \amp s \in x \]
\item Аксиома степени: существует $\mathcal{P}(x)$. Каково бы ни было множество
  $x$, существует множество, содержащее в точности все возможные подмножества
  множества $x$.
  \[ \forall x.\exists p.\forall y.\,y \in p \leftrightarrow y \subseteq x \]
\item Схема аксиом выделения: существует $\{t \in x \mid \phi(t)\}$. Для любого
  множества $x$ и любой формулы от одного аргумента $\phi(y)$ ($b$ не входит
  свободно в $\phi$), найдется $b$, в которое входят те и только те элементы из
  множества $x$, что $\phi(y)$ истинно.
  \[ \forall x.\exists b.\forall y.\,y \in b \leftrightarrow (y \in x \amp \phi(y)) \]
\end{itemize}
\end{dfn}
\begin{dfn}
Упорядоченной парой двух множеств $a$ и $b$ назовём $\{\{a\}, \{a, b\}\}$, или
$\langle a, b\rangle$.
\end{dfn}
\begin{thm}
Упорядоченную пару можно построить для любых множеств.
\end{thm}
\begin{proof}
Применить аксиому пары, теорему о существовании $\{X\}$, аксиому пары.
\end{proof}
\begin{dfn}[Типы порядка]
\begin{itemize}
\item Частичный порядок~--- рефлексивное, антисимметричное, транзитивное
  отношение.
\item Линейный порядок~--- частичный порядок, в котором любые два элемента
  сравнимы.
\item Полный (тотальный) порядок~--- обычно синоним линейного порядка; иногда
  требует, чтобы каждое подмножество имело наименьший элемент (более сильное
  условие).
\end{itemize}
\end{dfn}
\begin{dfn}[Аксиома бесконечности]
Инкремент: $x' \equiv x \cup \{x\}$.

Аксиома бесконечности. Существует $N$: $\varnothing \in N \amp \forall x.\,x \in
N \to x' \in N$. В $N$ есть всевозможные множества вида $\varnothing$,
$\{\varnothing\}$, $\{\varnothing, \{\varnothing\}\}$, $\{\varnothing,
\{\varnothing\}, \{\varnothing, \{\varnothing\}\}\}$, \dots
\end{dfn}
\begin{dfn}[Ординальные числа]
\begin{itemize}
\item Транзитивное множество $X$: $\forall x.\forall y.\,x \in y \amp y \in X
  \to x \in X$.
\item Ординал (порядковое число)~--- вполне упорядоченное отношением $(\in)$
  транзитивное множество.
\end{itemize}
\end{dfn}
\begin{thm}
Если $x, y$~--- ординалы, то $x = y$, или $x \in y$, или $y \in x$.
\end{thm}
$\omega$~--- наименьший предельный ординал.
\begin{thm}
$\omega$ существует.
\end{thm}
\begin{proof}
Пусть $\omega = \{x \in N \mid x \text{ конечен}\}$. Тогда:
\begin{itemize}
\item меньше $\omega$ предельных нет: если $\theta$ таков, что $\theta \in
  \omega$, тогда $\theta$ конечен.
\item $\omega$ предельный: Пусть $\theta$ таков, что $\theta' = \omega$. Тогда
  $\theta$ конечен и $\theta'$ тоже конечен.
\end{itemize}
\end{proof}
$\omega'$~--- тоже ординал.

\begin{dfn}
Порядковый тип множества~--- некоторое свойство, общее для всех множеств,
изоморфных относительно биективных отображений, сохраняющих порядок.

Порядковый тип вполне упорядоченного множества $\langle S, (\preceq)\rangle$~---
ординал $A$, для которого есть биективное отображение $f : S \to A$,
сохраняющее порядок: $a \preceq b$ тогда и только тогда, когда $f(a) \le f(b)$.

Множество $\mathbb{Z}$ не имеет порядкового типа (в смысле определения через
ординалы): оно не вполне упорядочено.
\end{dfn}
\begin{dfn}[Операции над ординалами]
\begin{itemize}
\item $a + b$~--- порядковый тип $a \uplus b$ (отмеченного объединения), причём
  $x_a < y_b$ при любых $x \in a$ и $y \in b$.
\item $a \cdot b$~--- порядковый тип $a \times b$, произведение упорядочено
  лексикографически: $\langle x_1, y_1\rangle < \langle x_2, y_2\rangle$, если
  $x_1 < x_2$ или $x_1 = x_2$ и $y_1 < y_2$.
\end{itemize}
\end{dfn}
\begin{dfn}
$\upb x$~--- верхняя грань множества ординалов, $\upb x = \bigcup_{a \in x} a$.
\[
\upb\{\varnothing', \varnothing'', \varnothing''''\} = \varnothing' \cup \varnothing'' \cup \varnothing'''' = \{\varnothing\} \cup \{\varnothing, \{\varnothing\}\} \cup \{\varnothing, \{\varnothing\}, \{\varnothing, \{\varnothing\}\}, \{\varnothing, \{\varnothing\}, \{\varnothing, \{\varnothing\}\}\}\} = \varnothing''''
\]
\end{dfn}
\begin{dfn}[Ординалы (порядковые числа) и порядок]
\begin{itemize}
\item Добавить элемент перед бесконечностью: $\mathbb{N}$ и $\mathbb{N}_0$.
  $1 + \omega = \omega$.
\item Добавить элемент после бесконечности ($+\infty$). $\omega + 1 \ne \omega$.
\end{itemize}
Упорядоченные пары натуральных чисел имеют порядковый тип $\omega^2$:
$\langle 3, 5\rangle < \langle 4, 3\rangle$, $\omega \cdot 3 + 5 < \omega \cdot
4 + 3$.

Списки натуральных чисел~--- порядковый тип $\omega^\omega$:
\[ \langle 3, 1, 4, 1, 5, 9\rangle : \omega^5 \cdot 3 + \omega^4 \cdot 1 + \omega^3 \cdot 4 + \omega^2 \cdot 1 + \omega^1 \cdot 5 + 9 \]
\end{dfn}

\section{Мощность множеств}
\begin{dfn}[Аксиома фундирования]
В каждом непустом множестве найдется элемент, не пересекающийся с исходным
множеством.
\[ \forall x.\,x = \varnothing \vee \exists y.\,y \in x \amp \forall z.\,z \in x \to z \notin y \]
Иными словами, в каждом множестве есть элемент, минимальный по отношению
$(\in)$.
\end{dfn}
Идея Рассела: каждому множеству припишем тип (тип пустого 0, тип множеств 1, тип
множеств множеств 2 и т.п.). Тогда конструкция невозможна: $\{x \mid x \in x\}$.
Аксиома фундирования позволяет определить функцию ранга:
\[ \rk(x) = \upb\{\rk(y) \mid y \in x\} \]
\begin{dfn}[Схема аксиом подстановки]
Пусть задана некоторая функция $f$, представимая в исчислении предикатов: то
есть задана некоторая формула $\varphi$, такая, что $f(x) = y$ тогда и только
тогда, когда $\varphi(x, y) \amp \exists!z.\,\varphi(x, z)$. Тогда для любого
множества $S$ существует множество $f(S)$~--- образ множества $S$ при
отображении $f$:
\[
\forall s.\,(\forall x.\forall y_1.\forall y_2.\,x \in s \amp \phi(x, y_1) \amp \phi(x, y_2) \to y_1 = y_2) \to (\exists t.\forall y.\,y \in t \leftrightarrow \exists x.\,x \in s \amp \phi(x, y))
\]
\end{dfn}
\begin{dfn}
Множество $A$ равномощно $B$ ($|A| = |B|$), если существует биекция $f : A \to
B$.
\end{dfn}
\begin{thm}[Кантора--Бернштейна]
Множество $A$ имеет мощность, не превышающую мощности $B$ ($|A| \le |B|$),
если существует инъекция $f : A \to B$. \emph{(Формулировка в исходном конспекте
неточна; сама теорема: если $|A| \le |B|$ и $|B| \le |A|$, то $|A| = |B|$.)}
\end{thm}
\begin{proof}
Избавимся от множества $B$: пусть $A_0 = A$; $A_1 = g(B)$; $A_{k+2} =
g(f(A_k))$. (В исходном конспекте здесь иллюстрация: убывающая цепочка $A_0
\supseteq A_1 \supseteq A_2 \supseteq \dots$ с пересечением $\bigcap A_k$ и
соответствующая цепочка $B_0 \supseteq B_1 \supseteq \dots$.)

Тогда, если существует $h : A_0 \to A_1$~--- биекция, то тогда $g^{-1} \circ h :
A \to B$~--- требуемая биекция.

Пусть $C_k = A_k \setminus A_{k+1}$. Тогда $g(f(C_k)) = g(f(A_k)) \setminus
g(f(A_{k+1})) = A_{k+2} \setminus A_{k+3} = C_{k+2}$.

Тогда определим $h(x)$ следующим образом:
\[
h(x) =
\begin{cases}
x, & x \in C_{2k+1} \vee x \in \bigcap A_k\\
g(f(x)), & x \in C_{2k}
\end{cases}
\]
\end{proof}
\begin{dfn}
Кардинальное число~--- наименьший ординал, не равномощный никакому меньшему:
$\forall x.\,x \in c \to |x| < |c|$.
\end{dfn}
\begin{thm}
Конечные ординалы~--- кардинальные числа.
\end{thm}
\begin{dfn}
Мощность множества ($|S|$)~--- равномощное ему кардинальное число.
\end{dfn}
\begin{thm}[Кантор]
$|\mathcal{P}(S)| > |S|$.
\end{thm}
\begin{proof}
Пусть $S = \{a, b, c, \dots\}$. (В исходном конспекте здесь рисунок~---
диагональная таблица.) Пусть $f : S \to \mathcal{P}(S)$~--- биекция. Тогда
$\sigma = \{y \in S \mid y \notin f(y)\}$. Пусть $f(x) = \sigma$. Но $x \in f(x)$
тогда и только тогда, когда $x \notin \sigma$, то есть $f(x) \ne \sigma$.
\end{proof}
\begin{dfn}
$\aleph_0 := |\omega|$; $\aleph_{k+1} := \min\{|a| \mid a$~--- ординал, $\aleph_k
< |a|\}$.
\end{dfn}
\begin{dfn}
$\beth_0 := |\omega|$; $\beth_{k+1} := |\mathcal{P}(\beth_k)|$.
\end{dfn}
Континуум-гипотеза (Г.~Кантор, 1877): $\aleph_1 = \beth_1$ (не существует
мощности, промежуточной между счётной и континуумом).

Обобщённая континуум-гипотеза: $\aleph_n = \beth_n$ при всех $n$.
\begin{dfn}
Утверждение $\alpha$ противоречит аксиоматике: $\proves \alpha$ ведёт к
противоречию.
\end{dfn}
\begin{dfn}
Утверждение $\alpha$ не зависит от аксиоматики: $\nvdash \alpha$ и $\nvdash
\neg\alpha$.
\end{dfn}
\begin{thm}
Утверждение $\aleph_1 = \beth_1$ не зависит от аксиоматики ZFC.
\end{thm}

\section{Модели, мощность}
\begin{dfn}
Пусть задана модель $\langle D, F_n, P_n\rangle$ для некоторой теории первого
порядка. Её мощностью будем считать мощность $D$.
\end{dfn}
\begin{dfn}
Пусть задана формальная теория с аксиомами $\{\alpha_n\}$. Её мощность~---
мощность множества $\{\alpha_n\}$. Формальная арифметика, исчисление предикатов,
исчисление высказываний~--- счётно-аксиоматизируемые.
\end{dfn}
\begin{dfn}
$\mathcal{M}' = \langle D', F'_n, P'_n\rangle$~--- элементарная подмодель
$\mathcal{M} = \langle D, F_n, P_n\rangle$, если:
\begin{itemize}
\item $D' \subseteq D$, $F'_n, P'_n$~--- сужение $F_n, P_n$ (замкнутое на $D'$).
\item $\mathcal{M} \models \varphi(x_1, \dots, x_n)$ тогда и только тогда, когда
  $\mathcal{M}' \models \varphi(x_1, \dots, x_n)$ при $x_i \in D'$.
\end{itemize}
\end{dfn}

\begin{thm}[Лёвенгейма--Сколема]
Пусть $T$~--- множество всех формул теории первого порядка. Пусть теория имеет
некоторую модель $\mathcal{M}$. Тогда найдётся элементарная подмодель
$\mathcal{M}'$, причём $|\mathcal{M}'| \le \max(\aleph_0, |T|)$.
\end{thm}
\begin{proof}
План:
\begin{enumerate}
\item Построим $D_0$~--- множество всех значений, которые упомянуты в языке
  теории.
\item Будем последовательно пополнять $D_i$: $D_0 \subseteq D_1 \subseteq D_2
  \dots$, следя за мощностью. $D' = \bigcup D_i$.
\item Покажем, что $\langle D', F_n, P_n\rangle$~--- требуемая подмодель.
\end{enumerate}
Пусть $\{f^0_k\}$~--- все 0-местные функциональные символы теории.
\begin{enumerate}
\item $D_0 = \{\llbracket f^0_k\rrbracket\}$, если есть хотя бы один $f^0_k$.
\item Если таких $f^0_k$ нет, возьмём какое-нибудь одно значение из $D$.
\end{enumerate}
Очевидно, $|D_0| \ge |T|$ \emph{(так в исходном конспекте)}.

Фиксируем некоторый $D_k$. Напомним, $T$~--- множество всех формул теории.
Рассмотрим $\varphi \in T$.
\begin{enumerate}
\item $\varphi$ не имеет свободных переменных~--- пропустим.
\item $\varphi$ имеет хотя бы одну свободную переменную $y$.
  \begin{enumerate}
  \item $\varphi(y, x_1, \dots, x_n)$ при $y, x_i \in D_k$ бывает истинным и
    ложным~--- ничего не меняем
  \item $\varphi(y, x_1, \dots, x_n)$ при $y \in D$ и $x_i \in D_k$ либо всегда
    истинен, либо всегда ложен~--- ничего не меняем
  \item $\varphi(y, x_1, \dots, x_n)$ при $y, x_i \in D_k$ тождественно истинен
    или ложен, но при $y' \in D \setminus D_k$ отличается~--- добавим $y'$ к
    $D_{k+1}$. Вместе добавим всевозможные $\llbracket\theta(y')\rrbracket$.
  \end{enumerate}
\end{enumerate}
\begin{enumerate}
\item Всего добавили не больше $|T| \cdot |T|$ (для каждой формулы $\varphi$,
  возможно, будет добавлен $y$~--- и всевозможные выражения $\theta(y)$,
  допустимые в языке), и $|D_0| \ge |T| \ge |T| \cdot |T|$, отсюда $|D_k| \ge
  |T| \cdot |T|$.
\item $|D'| = |\bigcup D_i| \ge |T| \cdot |T| \cdot \aleph_0$.
\item Тогда $|T| \cdot |T| \cdot \aleph_0 = \max(|T|, \aleph_0)$. Разберём
  случаи:
  \begin{enumerate}
  \item Если $|T| < \aleph_0$, тогда $(|T| \cdot |T|) \cdot \aleph_0 = \aleph_0$.
  \item Если $|T| \le \aleph_0$, тогда $(|T| \cdot |T|) \cdot \aleph_0 = |T|
    \cdot \aleph_0 = |T|$.
  \end{enumerate}
\item Итого, $|D'| \ge \max(|T|, \aleph_0)$.
\end{enumerate}
\emph{(Знаки неравенства в п.\ 1--4 приведены как в исходном конспекте; по
смыслу~--- оценки сверху, $\le$.)}

Докажем, что $\mathcal{M}'$~--- элементарная подмодель. Индукцией по структуре
формул $\tau \in T$ покажем, что все формулы можно вычислить, и что
$\llbracket\varphi\rrbracket_{\mathcal{M}'} = \llbracket\varphi\rrbracket_{\mathcal{M}}$.
\begin{enumerate}
\item База, 0 связок. $\tau \equiv P(f_1(x_1, \dots, x_n), \dots, f_{n}(x_1,
  \dots, x_n))$. Если $x_i \in D'$, то значит, добавлены на некоторых шагах
  (максимальный пусть $t$). Поэтому в $D_{t+1}$ можно вычислить формулу, и её
  значение сохранилось.
\item Переход. Пусть формулы из $k$ связок сохраняют значения. Рассмотрим
  $\tau$ с $k+1$ связкой.
  \begin{enumerate}
  \item $\tau \equiv \rho \star \sigma$~--- очевидно.
  \item $\tau \equiv \forall y.\varphi(y, x_1, \dots, x_n)$. Каждый $x_i$ добавлен
    на каком-то шаге~--- максимум $t$. Если $\varphi(y, x_1, \dots, x_n)$
    бывает истинен и ложен при $y_t, y_f \in D$, то $y_t, y_f \in D_{t+1}$ (по
    построению). Поэтому, если $\mathcal{M} \not\models \forall y.\varphi(y,
    x_1, \dots, x_n)$, то и $\mathcal{M}' \not\models \forall y.\varphi(y, x_1,
    \dots, x_n)$. Если же $\varphi(y, x_1, \dots, x_n)$ не меняется от $y$, то
    тем более $\llbracket\varphi\rrbracket_{\mathcal{M}'} =
    \llbracket\varphi\rrbracket_{\mathcal{M}}$.
  \item $\tau \equiv \exists y.\varphi(y, x_1, \dots, x_n)$~--- аналогично.
  \end{enumerate}
\end{enumerate}
\end{proof}

\begin{dfn}[«Парадокс» Сколема]
Как известно, $|\mathbb{R}| = |\mathcal{P}(\mathbb{N})| > |\mathbb{N}| =
\aleph_0$. Однако, ZFC~--- теория со счётным количеством формул. Значит,
существует счётная модель ZFC, то есть $|\mathbb{R}| = \aleph_0$. В чём ошибка?
У равенств разный смысл, первое~--- в предметном языке, второе~--- в метаязыке.
\end{dfn}

\section{Аксиома выбора}
\begin{axm}[выбора]
Из любого семейства дизъюнктных непустых множеств $\mathcal{A}$ можно выбрать
непустую трансверсаль~--- множество $S$, что $|S \cap A| = 1$ для каждого $A \in
\mathcal{A}$. Иначе, $S \in \times\mathcal{A}$.
\end{axm}

\begin{thm}[Функциональный вариант аксиомы выбора]
Пусть $\mathcal{A}$~--- семейство непустых множеств. Тогда существует $f :
\mathcal{A} \to \bigcup\mathcal{A}$, причём $\forall a.\,a \in \mathcal{A} \to
f(a) \in a$.
\end{thm}
\begin{proof}
Пусть $X(A) = \{\langle A, a\rangle \mid a \in A\}$, по семейству $\mathcal{A}$
рассмотрим $\{X(A) \mid A \in \mathcal{A}\}$
\begin{itemize}
\item непустых: если $A \in \mathcal{A}$, $A \ne \varnothing$, то $X(A) \ne
  \varnothing$;
\item дизъюнктное: если $A_0, A_1 \in \mathcal{A}$, $A_0 \ne A_1$, то $X(A_0)
  \cap X(A_1) = \varnothing$
\end{itemize}
тогда по аксиоме выбора $\exists f.\,f \in \times\mathcal{A}$.
\end{proof}

\begin{thm}[Лемма Цорна]
Если задано $\langle M, (\preceq)\rangle$ и для всякого линейно упорядоченного
$S \subseteq M$ выполнено $\upb_M S \ne \varnothing$, то в $M$ существует
максимальный элемент.
\end{thm}
\begin{thm}[Цермело]
На любом множестве можно задать полный порядок.
\end{thm}
\begin{thm}
У любой сюръективной функции существует частичная обратная.
\end{thm}
\begin{thm}
Аксиома выбора $\to$ лемма Цорна: без доказательства.
\end{thm}

Назовём (для данного раздела) упорядоченным множеством пару $\langle S,
(\prec_S)\rangle$. Отношение порядка $(\prec_S)$ может быть как строгим, так и
нестрогим.

\begin{dfn}
Будем говорить, что $\langle S, (\prec_S)\rangle$~--- начальный отрезок $\langle
T, (\prec_T)\rangle$, если:
\begin{itemize}
\item $S \subseteq T$;
\item если $a, b \in S$, то $a \prec_S b$ тогда и только тогда, когда $a \prec_T
  b$;
\item если $a \in S$, $b \in T \setminus S$, то $a \prec_T b$.
\end{itemize}
Будем обозначать это как $\langle S, (\prec_S)\rangle \sqsubseteq \langle T,
(\prec_T)\rangle$ или как $S \sqsubseteq T$, если порядок на множествах понятен
из контекста.
\end{dfn}
\begin{thm}
Отношение «быть начальным отрезком» является отношением нестрогого порядка.
\end{thm}
\begin{thm}[о верхней грани]
Если семейство упорядоченных множеств $X$ линейно упорядочено отношением «быть
начальным отрезком», то у него есть верхняя грань.
\end{thm}
\begin{proof}
Пусть $M = \bigcup\{T \mid \langle T, (\prec_T)\rangle \in X\}$ и $(\prec_M) =
\bigcup\{(\prec_T) \mid \langle T, (\prec_T)\rangle \in X\}$.

Покажем, что если $\langle A, (\prec_A)\rangle \in X$, то $A \sqsubseteq M$.
Рассмотрим определение:
\begin{itemize}
\item $A \subseteq M$~--- выполнено по построению $M$;
\item если $a, b \in A$, то $a \prec_A b$ влечёт $a \prec_M b$ (по построению
  $M$). Если же $a \prec_M b$, но $a \nprec_A b$, то существует $A'$, что $a, b
  \in A'$ и $a \prec_{A'} b$. Тогда $A \not\sqsubseteq A'$ и $A' \not\sqsubseteq
  A$, что невозможно по линейности порядка;
\item если $a \in A$, $b \in M \setminus A$, то найдётся $B$, что $b \in B$,
  отчего $a \prec_B b$ (так как $A \sqsubseteq B$) и $a \prec_M b$ (по
  построению $M$).
\end{itemize}
Тогда $\langle M, (\prec_M)\rangle$~--- требуемая верхняя грань.
\end{proof}

\begin{thm}
Лемма Цорна $\to$ теорема Цермело.
\end{thm}
\begin{proof}
Пусть $S = \{\langle P, (\prec)\rangle \mid P \subseteq X,\ (\prec)$ --- полный
порядок$\}$. Например, для $X = \{0, 1\}$ множество $S = \{\langle\varnothing,
\varnothing\rangle, \langle\{0\}, \varnothing\rangle, \langle\{1\},
\varnothing\rangle, \langle X, 0 \prec 1\rangle, \langle X, 1 \prec
0\rangle\}$. Введём порядок на $S$ как $(\sqsubseteq)$. Заметим, что это~---
частичный, но не линейный порядок. Например, $\langle X, 0 \prec 1\rangle$
несравним с $\langle X, 1 \prec 0\rangle$.

По теореме о верхней грани любое линейно упорядоченное подмножество $\langle T,
(\sqsubseteq)\rangle$ (где $T \subseteq S$) имеет верхнюю грань. Например, для
$\{\langle\varnothing, \varnothing\rangle, \langle\{0\}, \varnothing\rangle,
\langle X, 0 \prec 1\rangle\}$ это $\langle X, 0 \prec 1\rangle$.

По лемме Цорна тогда есть $\langle R, (\prec_R)\rangle = \max S$. Заметим, что $R
= X$, потому что иначе пусть $a \in X \setminus R$. Тогда положив $M = \langle R
\cup \{a\}, (\prec_R) \cup \{x \prec a \mid x \in R\}\rangle$ получим, что $M$
тоже вполне упорядоченное (и потому $M \in S$), значит, $R$ не максимальное.
\end{proof}

\begin{thm}
Теорема Цермело $\to$ у сюръективных функций существует частичная обратная.
\end{thm}
\begin{proof}
Рассмотрим сюръективную $f : A \to B$. Рассмотрим семейство $R_b = \{a \in A
\mid f(a) = b\}$. Построим полный порядок на каждом из $R_b$. Тогда $f^{-1}(b)
= \min R_b$.
\end{proof}
\begin{thm}
Существует частичная обратная у сюръективных функций $\to$ существует
трансверсаль у семейства непустых дизъюнктных множеств.
\end{thm}
\begin{proof}
Пусть дано семейство непустых дизъюнктных множеств $\mathcal{A}$. Рассмотрим $f
: \bigcup\mathcal{A} \to \mathcal{A}$, что $f(a) = \bigcup\{A \in \mathcal{A}
\mid a \in A\}$. Поскольку элементы $\mathcal{A}$ дизъюнктны, $f(a) \in
\mathcal{A}$ при всех $a$. Тогда существует $f^{-1} : \mathcal{A} \to
\bigcup\mathcal{A}$. Тогда $\{f^{-1}(A) \mid A \in \mathcal{A}\} \in
\times\mathcal{A}$.
\end{proof}

\section{Применение аксиомы выбора}
\begin{thm}[Диаконеску]
Если рассмотреть ИИП с ZFC, то для любого $P$ выполнено $\proves P \vee \neg P$.
\end{thm}
\begin{proof}
Рассмотрим $\mathcal{B} = \{0, 1\}$, $A_0 = \{x \in \mathcal{B} \mid x = 0 \vee
P\}$ и $A_1 = \{x \in \mathcal{B} \mid x = 1 \vee P\}$. $\{A_0, A_1\}$~---
семейство непустых множеств, и по акс.\ выбора существует $f : \{A_0, A_1\} \to
\bigcup A_i$, что $f(A_i) \in A_i$. (Если $P$, то $A_0 = A_1 = \{\mathcal{B}\}$.)
\begin{align*}
&\proves f(A_0) \in A_0 \amp f(A_1) \in A_1 &&\text{а.выбора: } f(A_i) \in A_i\\
&\proves f(A_0) \in \mathcal{B} \amp (f(A_0) = 0 \vee P) \amp f(A_1) \in \mathcal{B} \amp (f(A_1) = 1 \vee P) &&\text{а.выделения}\\
&\proves (f(A_0) = 0 \amp f(A_1) = 1) \vee P &&\text{Удал.\ ($\amp$) + дистр.}\\
&\proves P \vee f(A_0) \ne f(A_1) &&\text{$0 \ne 1$ и транз.}\\
&\proves P \to A_0 = A_1 &&\text{Определение } A_i\\
&\proves A_0 = A_1 \to f(A_0) = f(A_1) &&\text{Конгруэнтность}\\
&\proves f(A_0) \ne f(A_1) \to \neg P &&\text{Контрапозиция}\\
&\proves P \vee \neg P &&\text{Подставили}
\end{align*}
\end{proof}

\begin{thm}[конечного выбора]
Если $X_1 \ne \varnothing, \dots, X_n \ne \varnothing$, $X_i \cap X_j =
\varnothing$ при $i \ne j$, то $\times\{X_1, \dots, X_n\} \ne \varnothing$.
\end{thm}
\begin{proof}
\begin{itemize}
\item База: $n = 1$. Тогда $\exists x_1.\,x_1 \in X_1$, поэтому $\exists
  x_1.\,\{x_1\} \in \times\{X_1\}$.
\item Переход:
  \[
  \exists v.\,v \in \times\{X_{1,n}\} \to \exists x_{n+1}.\,x_{n+1} \in X_{n+1} \to v \cup \{x_{n+1}\} \in \times(X_{1,n} \cup \{X_{n+1}\})
  \]
\end{itemize}
\end{proof}
\begin{axm}[счётного выбора]
Для счётного семейства непустых множеств существует функция, каждому из которых
сопоставляющая один из своих элементов.
\end{axm}
\begin{axm}[зависимого выбора]
Если $\forall x \in E.\exists y \in E.\,x R y$, то существует последовательность
$x_n$: $\forall n.\,x_n R x_{n+1}$.
\end{axm}
\begin{dfn}
Наследственным свойством множества назовём такое свойство, которым обладает как
само множество, так и все его подмножества.
\end{dfn}
\begin{dfn}
Фундированным множеством назовём такое, которое не пересекается хотя бы с одним
своим элементом.
\end{dfn}
\begin{dfn}
Универсум фон Неймана $V$~--- все наследственные фундированные множества. При
наличии аксиомы фундирования можно показать, что $V = \bigcup_a V_a$, где:
\[
V_a =
\begin{cases}
\text{ничего}, & a = 0\\
\mathcal{P}(V_b), & a = b'\\
\bigcup_{b < a}(V_b), & a \text{ --- предельный}
\end{cases}
\]
\end{dfn}
\begin{dfn}
Конструктивный универсум $L = \bigcup_a L_a$, где:
\[
L_a =
\begin{cases}
\text{ничего}, & a = 0\\
\{\{x \in L_b \mid \varphi(x, t_1, \dots, t_k)\} \mid \varphi \text{ --- формула},\ t_i \in L_b\}, & a = b'\\
\bigcup_{b < a}(L_b), & a \text{ --- предельный}
\end{cases}
\]
\end{dfn}
\begin{dfn}[Аксиома конструктивности]
$V = L$, то есть допустимы только те фундированные множества, которые задаются
формулами.
\end{dfn}

\section{Индукция}
\begin{dfn}[Принцип математической индукции]
Какое бы ни было $\varphi(x)$, если $\varphi(0)$ и при всех $x$ выполнено
$\varphi(x) \to \varphi(x')$, то при всех $x$ выполнено и само $\varphi(x)$.
\end{dfn}
\begin{dfn}[Принцип полной математической индукции]
Какое бы ни было $\psi(x)$, если $\psi(0)$ и при всех $x$ выполнено $(\forall
t.\,t \le x \to \psi(t)) \to \psi(x')$, то при всех $x$ выполнено и само
$\psi(x)$.
\end{dfn}
\begin{dfn}
Назовём вполне упорядоченное отношением $(\in)$ множество $S$ наследственным
подмножеством $A$, если $\forall x.\,x \in A \to (\forall t.\,t \in x \to t \in
S) \to x \in S$.
\end{dfn}
\begin{thm}
Единственным наследственным подмножеством вполне упорядоченного множества
является оно само.
\end{thm}
\begin{proof}
Пусть $B \subseteq A$~--- наследственное и $B \ne A$. Тогда существует $a =
\min(A \setminus B)$. Тогда $(\forall t.\,t \in a \to t \in B) \to a \in B$ по
наследственности $B$, и выполнено $\forall t.\,t \in a \to t \in B$ (по
минимальности $a$). Значит, $a \in B$.
\end{proof}

\begin{thm}[Ограниченная трансфинитная индукция]
Если для $\varphi(x)$ (некоторого утверждения теории множеств) и некоторого
ординала $\varepsilon$ (ограничения) выполнено $\forall x.\,x \in \varepsilon
\to (\forall t.\,t \in x \to \varphi(t)) \to \varphi(x)$, то $\forall x.\,x \in
\varepsilon \to \varphi(x)$.
\end{thm}
\begin{proof}
Рассмотрим $S = \{x \in \varepsilon \mid \varphi(x)\}$. Тогда $x \in S$
равносильно $x \in \varepsilon \amp \varphi(x)$. Тогда перепишем: $\forall e.\,e
\in \varepsilon \to (\forall x.\,x \in e \to x \in S) \to e \in S$. Отсюда по
теореме о наследственных множествах $S = \varepsilon$.
\end{proof}

\begin{thm}[Неограниченная трансфинитная индукция]
Если для $\varphi(x)$ (некоторого утверждения теории множеств) выполнено
$\forall x.\,\text{ординал}(x) \to (\forall t.\,t \in x \to \varphi(t)) \to
\varphi(x)$, то $\forall x.\,\text{ординал}(x) \to \varphi(x)$.

ИЛИ

Для ординала $\varepsilon$ подмножество $S \subseteq \varepsilon$~---
наследственное, если и только если одновременно:
\begin{itemize}
\item Если $x \in \varepsilon$ и $x = \varnothing$, то $x \in S$;
\item Если $x \in \varepsilon$ и существует $y$: $y' = x$, то $y \in S \to x \in
  S$;
\item Если $x \in \varepsilon$ и $x$~--- предельный, то $(\forall t.\,t \in x \to
  t \in S) \to (x \in S)$.
\end{itemize}
\end{thm}
\begin{proof}
($\to$) очевидно. Докажем ($\leftarrow$): пусть $S$ не наследственное: $E :=
\{e \in \varepsilon \mid (\forall t.\,t \in e \to t \in S) \amp e \notin S\}$ и
$E \ne \varnothing$. Тогда пусть $e = \min E$.
\begin{enumerate}
\item $e = \varnothing$ или предельный. Тогда $(\forall t.\,t \in e \to t \in S)
  \to (e \in S)$.
\item $e = y'$. Тогда $y \in \varepsilon$ ($\varepsilon$~--- ординал) и $(\forall
  t.\,t \in y \to t \in S) \to (y \in S)$ (так как $e$ минимальный, для
  которого $S$ не наследственное). По условию, $(y \in S) \to (e \in S)$, отсюда
  $(\forall t.\,t \in e \to t \in S) \to (e \in S)$.
\end{enumerate}
\end{proof}

\begin{thm}
Если $\alpha$~--- кардинальное число, $\alpha \ge \aleph_0$, то $\alpha \cdot
\alpha = \alpha$.
\end{thm}
\begin{proof}
Трансфинитная индукция: $\varphi(x) := x < \omega \vee x \cdot x = x$.
\begin{enumerate}
\item База: $x = \varnothing$. Тогда $\varphi(\varnothing) \equiv \varnothing <
  \omega \vee |\varnothing \times \varnothing| = \varnothing$, что доказуемо.
\item Переход: $\forall y.\,y < x \to \varphi(y)$, тогда $\varphi(x)$. Три
  случая:
  \begin{enumerate}
  \item $x < \omega$. Тогда $\varphi(x)$ истинно (аналогично базе).
  \item $x = \omega$. Счётный случай (рассмотрим отдельно).
  \item $x > \omega$. Общий случай (рассмотрим отдельно).
  \end{enumerate}
\end{enumerate}

\emph{Счётный случай:} $\omega < \omega \vee |\omega \cdot \omega| = \omega$.
Тогда $\omega \times \omega$ упорядочим так: $\langle p, q\rangle \prec \langle
s, t\rangle$, если:
\begin{enumerate}
\item $\max(p, q) < \max(s, t)$
\item $\max(p, q) = \max(s, t)$ и $q < t$
\item $\max(p, q) = \max(s, t)$, $q = t$ и $p < s$
\end{enumerate}
Очевидно, можно построить биекцию между так упорядоченными значениями и
$\omega$.

\emph{Общий случай:} $|\alpha \cdot \alpha| = \alpha$. \todo{} Аналогично
счётному случаю, $\alpha \times \alpha$ упорядочим так: $\langle p, q\rangle
\prec \langle s, t\rangle$, если:
\begin{enumerate}
\item $p \cup q < s \cup t$
\item $p \cup q = s \cup t$ и $q < t$
\item $p \cup q = s \cup t$, $q = t$ и $p < s$
\end{enumerate}
\begin{itemize}
\item Легко заметить, что это~--- линейный порядок (показав, что $p \nprec q$ и
  $q \nprec p$ влечёт $p = q$)
\item \dots и полный порядок. Найти наименьший в $S \ne \varnothing$ возможно,
  рассмотрев $m_1 := \min\{p \cup q \mid \langle p, q\rangle \in S\}$ и $M_1 :=
  \{\langle p, q\rangle \mid \langle p, q\rangle \in S,\ p \cup q = m_1\}$,
  затем $m_2 := \min\{q \mid \langle p, q\rangle \in M_1\}$, $M_2 := \{\langle p,
  q\rangle \mid \langle p, q\rangle \in M_1,\ q = m_2\}$. Тогда требуемым
  наименьшим в $S$ будет $\min\{p \mid \langle p, q\rangle \in M_2\}$
\item Тогда $\langle\alpha \times \alpha, (\prec)\rangle$ соответствует какой-то
  ординал $\tau$ и сохраняющая порядок биекция $t : \tau \to \alpha \times
  \alpha$.
\item Заметим, что $x < \omega$ тогда и только тогда, когда $\bigcup(\bigcup
  t(x)) < \omega$ (очевидно из того, что $|\{z \mid \text{ординал}(z),\ z <
  x\}| = |\{p \mid p \prec t(x)\}|$).
\item Покажем, что $|\tau| = \alpha$. Докажем $\tau = \alpha$. Очевидно, что
  $\tau \ge \alpha$ (так как $|\tau| = |\alpha \times \alpha| \ge \alpha$). Но
  пусть $\tau > \alpha$.
\item Тогда $t(\alpha) = \langle\zeta, \eta\rangle$ определено (у $\alpha$ есть
  образ).
\item Пусть $\sigma := \zeta \cup \eta$. Очевидно, $\langle\zeta, \eta\rangle
  \preceq \langle\sigma, \sigma\rangle$ и $\sigma \in \alpha$.
\item Каков образ $t$ на этом начальном отрезке? $\{t(x) \mid x < \alpha\}
  \subseteq \{\langle p, q\rangle \mid p, q < \sigma\}$. Поэтому $\alpha <
  |(\sigma+1) \times (\sigma+1)|$.
\item С другой стороны, $\sigma < \alpha$. Поскольку $\alpha$~--- кардинал (т.е.,
  в частности, предельный ординал), то $\sigma + 1 < \alpha$ и $|\sigma + 1| <
  \alpha$.
\item По предположению индукции, $|\sigma + 1| < \omega \vee |\sigma + 1| =
  |\sigma + 1| \cdot |\sigma + 1|$, по свойствам $(\prec)$ имеем $\sigma >
  \omega$.
\item Отсюда $\alpha < |(\sigma+1) \times (\sigma+1)| = |\sigma + 1| < \alpha$,
  что невозможно.
\end{itemize}
\end{proof}

\section{Система $S_\infty$}
\begin{dfn}
Введем исчисление $S_\infty$.
\begin{enumerate}
\item Язык: связки $\neg$, $\vee$, $\forall$, $=$; нелогические символы:
  $(+)$, $(\cdot)$, $(')$, $0$; переменные: $x$.
\item Аксиомы: все истинные формулы вида $\theta_1 = \theta_2$; все истинные
  отрицания формул вида $\neg\theta_1 = \theta_2$ ($\theta_i$~--- термы без
  переменных).
\item Структурные (слабые) правила:
  \[
  \infer{\zeta \vee \beta \vee \alpha \vee \delta}{\zeta \vee \alpha \vee \beta \vee \delta}
  \qquad
  \infer{\alpha \vee \delta}{\alpha \vee \alpha \vee \delta}
  \]
  Сильные правила:
  \[
  \infer{\alpha \vee \beta}{\beta}
  \qquad
  \infer{\neg(\alpha \vee \beta) \vee \delta}{\neg\alpha \vee \delta & \neg\beta \vee \delta}
  \qquad
  \infer{\neg\neg\alpha \vee \delta}{\alpha \vee \delta}
  \qquad
  \infer{(\neg\forall x.\alpha) \vee \delta}{\neg\alpha[x := \theta] \vee \delta}
  \]
  Формулы в правилах, обозначенные буквами $\zeta$ и $\delta$, называются
  боковыми и могут отсутствовать.
\item Бесконечная индукция:
  \[
  \infer{(\forall x.\alpha) \vee \delta}{\alpha[x := \overline{0}] \vee \delta & \alpha[x := \overline{1}] \vee \delta & \alpha[x := \overline{2}] \vee \delta & \dots}
  \]
\item Сечение:
  \[
  \infer{\zeta \vee \delta}{\zeta \vee \alpha & \neg\alpha \vee \delta}
  \]
  Здесь $\alpha$~--- секущая формула, число связок в $\neg\alpha$~--- степень
  сечения. В отличие от других правил, в правиле сечения хотя бы одна из боковых
  формул $\zeta$ или $\delta$ должна присутствовать.
\end{enumerate}
\begin{enumerate}
\item Доказательства образуют деревья.
\item Каждой формуле в дереве сопоставим порядковое число (ординал).
\item Порядковое число заключения любого неструктурного правила строго больше
  порядкового числа его посылок (больше или равно в случае структурного
  правила).
  \[
  \infer{(\neg\neg\forall x.\neg x' = 0)_{\omega+1}}{
    \infer{(\forall x.\neg x' = 0)_\omega}{(\neg 1 = 0)_1 & (\neg 2 = 0)_2 & (\neg 3 = 0)_4 & (\neg 4 = 0)_8 & \dots}}
  \]
\item Существует конечная максимальная степень сечения в дереве (назовём её
  степенью вывода).
\end{enumerate}
\end{dfn}

\begin{thm}
Если $\alpha$ имеет вывод степени $m > 0$ порядка $t$, то можно найти вывод
степени строго меньшей $m$ с порядком $2^t$.
\end{thm}
\begin{proof}
Трансфинитная индукция. Пусть для всех деревьев порядка $t_1 < t$ условие
выполнено. Покажем, что оно выполнено для порядка $t$. Рассмотрим заключительное
правило. Это может быть\dots
\begin{enumerate}
\item Не сечение.
\item Сечение, секущая формула~--- элементарная.
\item Сечение, секущая формула~--- $\neg\alpha$.
\item Сечение, секущая формула~--- $\alpha \vee \beta$.
\item Сечение, секущая формула~--- $\forall x.\alpha$.
\end{enumerate}

\emph{Случай 1. Не сечение.}
\[
\infer{(\alpha)_t}{(\pi_0)_{t_0} & (\pi_1)_{t_1} & (\pi_2)_{t_2} & \dots}
\]
Заменим доказательства посылок $(\pi_i)_{t_i}$ на $(\pi'_i)_{2^{t_i}}$ по
индукционному предположению.
\begin{enumerate}
\item Поскольку степени посылок $m'_i < m_i$, то $\max m'_i < \max m_i$.
\item Поскольку $t_i \le t$, то $2^{t_i} \le 2^t$.
\end{enumerate}

\emph{Случай 5. Сечение с формулой вида $\forall x.\alpha$.}
\[
\infer{\zeta \vee \delta}{\zeta \vee \forall x.\alpha & (\neg\forall x.\alpha) \vee \delta}
\]
Причём степень и порядок выводов компонент, соответственно, $(m_1, t_1)$ и
$(m_2, t_2)$.
\begin{enumerate}
\item По индукции, вывод $\zeta \vee \forall x.\alpha$ можно упростить до
  $(m'_1, 2^{t_1})$.
\item По обратимости, можно построить вывод $\zeta \vee \alpha[x := \theta]$ за
  $(m'_1, 2^{t_1})$.
\item В формуле $(\neg\forall x.\alpha) \vee \delta$ формула $\neg\forall
  x.\alpha$ получена либо ослаблением, либо квантификацией из
  $\neg\alpha[x := \theta_k] \vee \delta_k$.
  \begin{enumerate}
  \item Каждое правило квантификации заменим на:
    \[
    \infer{\zeta \vee \delta_k}{\zeta \vee \alpha[x := \theta_k] & (\neg\alpha[x := \theta_k]) \vee \delta_k}
    \]
  \item Остальные вхождения $\neg\forall x.\alpha$ заменим на $\zeta$ (в
    правилах ослабления).
  \end{enumerate}
\item В получившемся дереве меньше степень~--- так как в $\neg\alpha[x :=
  \theta]$ меньше связок, чем в $\neg\forall x.\alpha$.
\end{enumerate}
\end{proof}

\begin{dfn}[Итерационная экспонента]
\[
(a\uparrow)^m(t) =
\begin{cases}
t, & m = 0\\
a^{(a\uparrow)^{m-1}(t)}, & m > 0
\end{cases}
\]
\end{dfn}
\begin{thm}[об устранении сечений]
Если $\proves_\infty \sigma$ степени $m$ порядка $t$, то найдётся доказательство
без сечений порядка $(2\uparrow)^m(t)$.
\end{thm}
\begin{dfn}
$\varepsilon_0$~--- неподвижная точка $\varepsilon_0 = \omega^{\varepsilon_0}$.
Иначе говоря, $\varepsilon_0 = \{\omega, \omega^\omega, \omega^{\omega^\omega},
(\omega\uparrow)^3(\omega), (\omega\uparrow)^4(\omega), \dots\}$.
\end{dfn}
Очевидно, что теорема об устранении сечений может быть доказана трансфинитной
индукцией до ординала $\varepsilon_0$ (максимальный порядок дерева вывода, при
правильной нумерации вершин).

\begin{lmm}
Если $\proves_\infty \alpha$ и $\proves_\infty \neg\alpha$, тогда $\proves_\infty
\neg 0 = 0$.
\end{lmm}
\begin{thm}[о непротиворечивости формальной арифметики]
$\nvdash_\infty \neg 0 = 0$.
\end{thm}
\begin{proof}
Пусть $\proves_\infty \neg 0 = 0$, устраним сечения и рассмотрим заключительное
правило.
\begin{enumerate}
\item Правило де Моргана? Нет отрицаний дизъюнкции ($\neg(\alpha \vee \beta)
  \vee \delta$).
\item Отрицание? Нет двойного отрицания ($\neg\neg\alpha \vee \delta$).
\item Бесконечная индукция или квантификация? Нет квантора.
\item Ослабление? Нет дизъюнкции ($\alpha \vee \beta$), хотя $\beta$ обязана
  присутствовать.
\item Сечение? Исключено по условию.
\end{enumerate}
То есть, неизбежно, $\neg 0 = 0$~--- аксиома, что также неверно.
\end{proof}

\section{Метод резолюции}
Сколемизация не сохраняет общезначимость.

Дана формула $\alpha$.
\begin{enumerate}
\item Упростим формулу~--- поверхностные кванторы всеобщности, сколемизация.
  Умеем строить формулу $\beta$: $\beta := \forall x_1.\forall x_2.\forall
  x_k.\,\delta_1(x_1, \dots, x_k) \amp \dots \amp \delta_n(x_1, \dots, x_k)$.

  $\alpha$ доказуема тогда и только тогда, когда при всех оценках предикатных и
  функциональных символов найдётся значение сколемовских функций $e_k$, при
  которых $\beta$ всегда истинна (слоёный пирог из кванторов).
\item Упрощаем предметное множество~--- заменили произвольный $D$ на эрбранов
  универсум $H$. Выполнимость формулы эквивалентна выполнимости на эрбрановом
  универсуме.
\item Осталось избавиться от кванторов всеобщности и организовать правильный
  перебор (эрбранов универсум может быть бесконечным).
\end{enumerate}
\begin{dfn}
Эрбранов универсум $H_\varphi$~--- всевозможные комбинации функциональных
символов из формулы $\varphi$. Если в формуле нет нульместных функциональных
символов, к множеству символов формулы добавляется свежий нульместный
функциональный символ $a$ и все комбинации с его участием.
\end{dfn}
\begin{dfn}
Если $\varphi$~--- бескванторная формула, то её эрбранова оценка задаётся как
$\langle H_\varphi, F, P, E\rangle$, функции $F$ определяются как текстовые
подстановки $\llbracket f(\theta)\rrbracket = \text{``}f(\text{''} +\!+
\llbracket\theta\rrbracket +\!+ \text{``)''}$, предикаты $P$ задаются
перечислением истинных.
\end{dfn}
\begin{dfn}
Система дизъюнктов $S = \{\delta_1, \dots, \delta_n\}$ противоречива, если для
каждой оценки $M = \langle D, P, F, E\rangle$ найдётся $\delta_t$ и такой набор
$\vec d \in D$, что $\llbracket\delta_t\rrbracket^{\vec x := \vec d} = \False$.
\end{dfn}
\begin{dfn}
Дизъюнкт с подставленными значениями из эрбранового универсума $H_\beta$ вместо
переменных называется основным примером формулы $\beta$.
\end{dfn}
\begin{dfn}
Система основных примеров~--- все основные примеры, опровергаемые хоть при
какой-то эрбрановой оценке $\mathcal{M}$:
\[
\mathcal{E}_S = \{\delta_t[\vec x := \vec d] \mid \text{существует } \mathcal{M},\ \text{что } \llbracket\delta_t[\vec x := \vec d]\rrbracket(M) = \False;\ d_i \in H_\beta\}
\]
\end{dfn}
\begin{thm}[Гёделя о компактности]
Если $\Gamma$~--- некоторое семейство бескванторных формул, то $\Gamma$ имеет
модель тогда и только тогда, когда любое его конечное подмножество имеет модель.
\end{thm}
\begin{thm}[Эрбрана]
Система дизъюнктов $S$ противоречива тогда и только тогда, когда у
$\mathcal{E}_S$ существует конечное противоречивое в эрбрановой интерпретации
подмножество.
\end{thm}
\begin{proof}
($\Leftarrow$) Пусть $\{\varepsilon_1, \dots, \varepsilon_t\} \subseteq
\mathcal{E}_S$ противоречиво, $\varepsilon_i = \delta_{m_i}[\vec x := \vec
d_i]$, где $\vec d_i$~--- набор значений из $H$. То есть, для любой эрбрановой
оценки $M$ существует $\varepsilon_p$, что $\llbracket\varepsilon_p\rrbracket_M =
\False$. Отсюда, по теореме о выполнимости $S$ тоже противоречива.

($\Rightarrow$) Если $S$ противоречива, то $\mathcal{E}_S$ противоречива. Тогда у
неё нет модели. Тогда у неё найдётся конечное противоречивое подмножество
(компактность). Возможно убедиться в невыполнимости за конечное время.
\end{proof}

\subsection*{Общая схема алгоритма}
Цель алгоритма: убедиться, что $\alpha$ доказуемо.
\begin{enumerate}
\item По формуле $\alpha$ строим её отрицание $\neg\alpha$.
\item Приводим к виду с поверхностными кванторами, проводим сколемизацию,
  находим КНФ: $\beta = \forall x_1.\dots.\forall x_k.\,\delta_1 \amp \dots \amp
  \delta_n$.
\item Убедимся, что при любом $D$ и значениях функциональных и предикатных
  символов и сколемовских функций $e_k$ найдутся $d_i \in D$, что один из
  дизъюнктов $\delta_t$ при подстановке $\vec x := \vec d$ ложный.
\item Для этого строим универсум Эрбрана $H$, и систему основных примеров
  $\mathcal{E}_S$, её противоречивость эквивалентна невыполнимости $\beta$.
\item Конечное противоречивое подмножество обязательно находится в каком-то
  начальном отрезке $\{\varepsilon_1, \dots, \varepsilon_t\} \subseteq
  \mathcal{E}_S$ (если оно есть).
\end{enumerate}

\begin{dfn}[Правило резолюции (исчисление высказываний)]
Пусть даны два дизъюнкта, $\alpha_1 \vee \beta$ и $\alpha_2 \vee \neg\beta$. Тогда
следующее правило вывода называется правилом резолюции:
\[ \infer{\alpha_1 \vee \alpha_2}{\alpha_1 \vee \beta & \alpha_2 \vee \neg\beta} \]
\end{dfn}
\begin{dfn}
Алгебраический терм $\theta := x \mid f(\theta_1, \dots, \theta_n)$, где
$x$~--- переменная, $f(\theta_1, \dots, \theta_n)$~--- применение функции.
Напомним, что константы~--- нульместные функциональные символы, собственно
переменные будем обозначать последними буквами латинского алфавита.
\end{dfn}
\begin{dfn}
Система уравнений в алгебраических термах $\begin{cases}\theta_1 = \sigma_1\\
\vdots\\ \theta_n = \sigma_n\end{cases}$, где $\theta_i$ и $\sigma_i$~--- термы.
\end{dfn}
\begin{dfn}
$\{x_i\} = X$~--- множество переменных, $\{\theta_i\} = T$~--- множество термов.
\end{dfn}
\begin{dfn}
Подстановка~--- отображение вида $\pi_0 : X \to T$, тождественное почти везде (за
исключением конечного числа переменных). $\pi_0(x)$ может быть либо $\pi_0(x) =
\theta_i$, либо $\pi_0(x) = x$. Доопределим $\pi : T \to T$, где
\begin{enumerate}
\item $\pi(x) = \pi_0(x)$
\item $\pi(f(\theta_1, \dots, \theta_k)) = f(\pi(\theta_1), \dots, \pi(\theta_k))$
\end{enumerate}
Решить уравнение в алгебраических термах~--- найти такую наиболее общую
подстановку $\pi$, что $\pi(\theta_1) = \pi(\theta_2)$. Наиболее общая
подстановка~--- такая, для которой другие подстановки являются её частными
случаями.
\end{dfn}
\begin{dfn}
Пусть даны формулы $\alpha$ и $\beta$. Тогда решением задачи унификации будет
такая наиболее общая подстановка $\pi = \mathcal{U}[\alpha, \beta]$, что
$\pi(\alpha) = \pi(\beta)$. Также, $\pi$ назовём наиболее общим унификатором.
\end{dfn}
\begin{dfn}[Правило резолюции (исчисление предикатов)]
Пусть $\sigma_1$ и $\sigma_2$~--- подстановки, заменяющие переменные в формуле на
свежие. Тогда правило резолюции выглядит так: \todo{}
\end{dfn}

\section{Информация о курсе}
Поток~--- y2024. Группы M3232--M3239. Преподаватель~--- Штукенберг Дмитрий
Григорьевич.

\end{document}
